\documentclass[11pt]{article}
\usepackage{tcee}
\addbibresource{3B19.bib}
\tcetema[La coordinación internacional de políticas económicas]{3.B.19}{La coordinación internacional de políticas económicas. Aspectos teóricos y prácticos. El G-20, la OCDE y otros foros internacionales.}
\tcefecha{11/11/2026}
\begin{document}
\tceetitulo

\section*{Introducción}
\minutos{2,5}

\begin{itemize}
\item \textbf{\emph{Enganche}}:
  \begin{itemize}
  \item «Cada uno en su casa y Dios en la de todos», dice el refrán.
  \item El 22 y el 23 de noviembre de 2025 se celebró en Johannesburgo la primera cumbre de líderes del G-20 en África. Estados Unidos, la primera economía del mundo y el país que preside el grupo en 2026, no acudió, y los demás aprobaron su declaración el primer día, sin él\nocite{G20Johannesburgo2025}.
  \item 16 años antes, en Londres, el mismo G-20 había acordado la mayor expansión fiscal concertada de la historia para sacar juntos a la economía mundial de la recesión\nocite{G20Londres2009}.
  \item Entre esas 2 cumbres está la pregunta de la exposición: \emph{¿basta con que cada país ponga orden en su casa, o hay problemas que solo se arreglan entre todos?}
  \end{itemize}

\item \textbf{\emph{Relevancia}}:
  \begin{itemize}
  \item \emph{Para la economía mundial}: con la integración comercial y financiera, la política de un país grande se nota fuera. La de la Reserva Federal mueve el crédito y los precios de los activos de todo el mundo\nocite{Rey2013Dilema,MirandaAgrippinoRey2020}, y como buena parte del comercio se factura en dólares, mueve también los precios del comercio entre terceros países\nocite{Gopinath2020DCP}.
  \item \emph{Para la política económica de hoy}: el FMI advierte de que los desequilibrios exteriores crecen y pide reequilibrar entre países, con más demanda interna en los que tienen superávit y consolidación en los deficitarios, y recuerda que el resultado es mejor si actúan todos a la vez\nocite{FMI2026ESR}. Es una petición de coordinación, en un momento de tensiones comerciales y de foros divididos.
  \item \emph{Para España}: es invitado permanente del G-20, miembro fundador de la OCDE y miembro del Consejo de Estabilidad Financiera\nocite{FSB2026Miembros}, y dentro del euro vive la forma más intensa de coordinación, la gobernanza económica de la Unión \vertema{3.B.45}. Su respuesta fiscal a la crisis de 2008-2009 se enmarcó en el estímulo concertado del G-20 \vertema{4.A.28}.
  \item \emph{Para la teoría}: el tema une la macroeconomía abierta \vertema{3.A.32}, la teoría de juegos \vertema{3.A.14} y la economía política internacional.
  \end{itemize}

\item \textbf{\emph{Contextualización}}:
  \begin{itemize}
  \item \emph{Las definiciones de partida}:
    \begin{itemize}
    \item la \concepto{cooperación} internacional es una escalera con varios peldaños: el más bajo es intercambiar información sobre la coyuntura y los planes de cada país (el trabajo diario de la OCDE o del FMI); el siguiente, ponerse de acuerdo en una acción puntual; el más alto, comprometerse con unas reglas permanentes;
    \item la \concepto{coordinación} es la forma más exigente de cooperación: cada país cambia su política, de mutuo acuerdo, para tener en cuenta los efectos que tiene sobre los demás. En términos de teoría de juegos, es pasar del equilibrio no cooperativo de \autor{Nash}\footnote{\autor{John Nash} recibió el Premio Nobel de Economía en 1994, con \autor{John Harsanyi} y \autor{Reinhard Selten}, «por su análisis pionero de los equilibrios en la teoría de los juegos no cooperativos»\nocite{Nobel1994Nash}.} a una solución negociada\nocite{Frankel2016Coordination};
    \item puede ser \emph{puntual} (el estímulo fiscal del G-20 en 2009) o \emph{permanente} (las paridades de Bretton Woods, el Pacto de Estabilidad y Crecimiento);
    \item puede tocar las políticas de \emph{demanda} (monetaria y fiscal) o las de \emph{oferta} (comercio, competencia, fiscalidad, clima). La exposición se centra en las de demanda; las de oferta salen con la reforma estructural de \autor{Corsetti} y \autor{Pesenti} y con la OCDE y el G-20.
    \end{itemize}
  \item \emph{Las ideas han ido y venido en 3 etapas}\nocite{Bergin2008Palgrave}:
    \begin{itemize}
    \item la primera aportación fue contraria a coordinar: \textcite{Friedman1953Case}\footnote{\autor{Milton Friedman} recibió el Premio Nobel de Economía en 1976 «por sus logros en el análisis del consumo, la historia y la teoría monetarias y por demostrar la complejidad de la política de estabilización»\nocite{Nobel1976Friedman}.} sostuvo que, con tipos de cambio flexibles, el mercado aísla a cada país de las perturbaciones de sus vecinos y cada uno puede ocuparse de lo suyo;
    \item en los 70 y los 80, los modelos keynesianos de 2 países y la teoría de juegos dieron argumentos a favor\nocite{Hamada1976Strategic,OudizSachs1984,CanzoneriGray1985}, y la práctica los siguió: Bonn en 1978, el Plaza en 1985, el Louvre en 1987\nocite{Frankel2016Coordination};
    \item desde finales de los 90, la nueva macroeconomía abierta, con microfundamentos, volvió al escepticismo: las ganancias de coordinar son pequeñas\nocite{ObstfeldRogoff2002Global}. La crisis de 2008 devolvió la coordinación a la práctica, con el G-20.
    \end{itemize}
  \item \emph{La convención y el ejemplo}:
    \begin{itemize}
    \item el tipo de cambio $E$ es directo (unidades de moneda nacional por unidad de moneda extranjera), de modo que una subida es una depreciación \vertema{3.B.13};
    \item la exposición sigue a 2 países grandes, A y B, los mismos en todos los modelos: A, la zona del euro, y B, Estados Unidos.
    \end{itemize}
  \end{itemize}

\item \textbf{\emph{Problemática}}:
  \begin{itemize}
  \item ¿Por qué coordinar? ¿Qué efectos tiene la política de un país sobre los demás, y por qué canales?
  \item ¿Cuándo es deseable coordinar? ¿Cuánto se gana, y por qué se ve tan poco?
  \item ¿Cómo se coordina en la práctica: con reglas o con acuerdos, y en qué foros?
  \end{itemize}

\item \textbf{\emph{Estructura}}:
  \begin{itemize}
  \item La exposición sigue el título, de la teoría a la práctica y de esta a los foros:
    \begin{itemize}
    \item \concepto{I. Análisis teórico} (\ref{sec:teoria}): cuándo es \emph{necesaria} la coordinación (hay efectos desbordamiento, con \autor{Mundell}\footnote{\autor{Robert Mundell} recibió el Premio Nobel de Economía en 1999 «por su análisis de la política monetaria y fiscal con distintos regímenes cambiarios y su análisis de las áreas monetarias óptimas»\nocite{Nobel1999Mundell}.}-\autor{Fleming} y con \autor{Corsetti} y \autor{Pesenti}) y cuándo es \emph{deseable} (el equilibrio sin coordinar es ineficiente, con el diagrama de \autor{Hamada});
    \item \concepto{II. La cooperación en la historia reciente} (\ref{sec:historia}): de Bonn y el Plaza a Londres 2009 y la guerra de divisas, por qué se coordina tan poco y de qué formas;
    \item \concepto{III. Economía política: los mecanismos} (\ref{sec:economia-politica}): por qué hacen falta foros y qué hacen el G-20, la OCDE y los demás.
    \end{itemize}
  \end{itemize}
\end{itemize}

\begin{esquema}[unbreakable]
\textbf{3.B.19. La coordinación internacional de políticas económicas}\\[2pt]
Introducción\esqtiempo{2,5}\\
I. Análisis teórico: efectos desbordamiento y deseabilidad de la coordinación\esqtiempo{12,5}\\
\hspace*{1.5em}1. Justificación. Condición necesaria: los efectos desbordamiento\\
\hspace*{1.5em}2. Deseabilidad. Condición suficiente: las ganancias. La teoría de juegos\\
II. La cooperación económica en la historia macroeconómica reciente\esqtiempo{6}\\
\hspace*{1.5em}1. Ejemplos de cooperación\\
\hspace*{1.5em}2. Obstáculos a la coordinación\\
\hspace*{1.5em}3. Formas de la coordinación\\
III. Economía política de la coordinación: mecanismos de implementación\esqtiempo{8}\\
\hspace*{1.5em}1. La problemática de teoría de juegos: por qué hacen falta foros\\
\hspace*{1.5em}2. Los foros de coordinación: el G-20, la OCDE y otros\\
Conclusión\esqtiempo{1}
\end{esquema}

\begin{notaopositor}
Reparto del cante: introducción, \tiempo{2,5}; I, \tiempo{12,5} (I.1, \tiempo{8,5}: los canales, \tiempo{0,5}; \autor{Mundell}-\autor{Fleming} con 2 países, \tiempo{3,5}; \autor{Corsetti} y \autor{Pesenti}, \tiempo{4,5}; I.2, \tiempo{4}: \autor{Hamada}, \tiempo{2,5}; la repetición, \tiempo{1}; evidencia y valoración, \tiempo{0,5}); II, \tiempo{6} (II.1, \tiempo{4}; II.2, \tiempo{1,5}; II.3, \tiempo{0,5}); III, \tiempo{8} (III.1, \tiempo{1,5}; III.2, \tiempo{6,5}), y conclusión, \tiempo{1}. Todo el tema vuelve a los mismos 2 países, la zona del euro (A) y Estados Unidos (B): en los paneles de \autor{Mundell}-\autor{Fleming} y de \autor{Corsetti} y \autor{Pesenti}, en los ejes de \autor{Hamada}, en la matriz de la locomotora y, al final, sentados en el G-20. Los gráficos y figuras imprescindibles son la expansión fiscal con 2 países, el de \autor{Corsetti} y \autor{Pesenti}, el diagrama de \autor{Hamada}, la línea del tiempo y la figura de la escalera de la cooperación; las funciones clave, las 3 curvas de \autor{Corsetti} y \autor{Pesenti} \eqref{eq:corsetti-pesenti}, la pérdida de cada país \eqref{eq:hamada} y la condición del gatillo \eqref{eq:gatillo}. Son de profundización (no se cantan): la política monetaria con 2 países (basta su fila de la tabla), el tipo de cambio fijo, el caso del dólar y la reforma estructural de \autor{Corsetti} y \autor{Pesenti} (anexo~\ref{anexo:corsetti-pesenti}), las tablas de los otros 2 juegos (anexo~\ref{anexo:juegos}), la amplitud, la profundidad y el tamaño de los foros en detalle, los actos jurídicos de la OCDE y el anexo~\ref{anexo:superavit}.
\end{notaopositor}

\begin{anotaciones}
\begin{itemize}
\item Para ampliar: la voz «International Policy Coordination» del \emph{New Palgrave Dictionary of Economics}\nocite{Bergin2008Palgrave}, el repaso de \textcite{Frankel2016Coordination}, con los juegos de la locomotora, la disciplina fiscal y la guerra de divisas, y la economía política de la coordinación monetaria de \textcite{BrozFrieden2013}.
\item El blog del FMI sobre el «monstruo del lago Ness»\nocite{BlanchardOstryGhosh2013} y su versión larga, con los obstáculos y cómo superarlos\nocite{OstryGhosh2013}.
\item El traspaso del tipo de cambio y el dólar: la entrevista a \autor{Gita Gopinath} «Analyze This! Exchange Rates and Trade»\nocite{FMIAnalyzeThisGopinath} y el capítulo 2 del Informe del Sector Exterior de 2019\nocite{FMI2019ESR}.
\end{itemize}
\end{anotaciones}


% ==========================================================================
\section{Análisis teórico:\cortetitulo efectos desbordamiento y deseabilidad de la coordinación}\label{sec:teoria}
\minutos{12,5}

\begin{itemize}
\item Empezamos por la teoría: qué tiene que pasar para que coordinar tenga sentido. Son 2 condiciones, una necesaria y otra suficiente.
\item \emph{Conclusión anticipada}:
  \begin{itemize}
  \item la condición necesaria es que haya \emph{efectos desbordamiento}: la política de un país grande mueve la renta y el bienestar de sus socios. Su signo depende del modelo: con \autor{Mundell}-\autor{Fleming}, la expansión fiscal tira del vecino y la monetaria lo empobrece; con \autor{Corsetti} y \autor{Pesenti}, lo decide la moneda en que se fijan los precios;
  \item la condición suficiente es que coordinar \emph{compense}: la teoría de juegos muestra que, con desbordamientos, actuar cada uno por su cuenta es ineficiente, pero la solución cooperativa no es un equilibrio y cada país tiene incentivos a romperla;
  \item los desbordamientos son condición necesaria, no suficiente: en algunos modelos los hay y, aun así, coordinar no mejora nada.
  \end{itemize}
\end{itemize}

% --------------------------------------------------------------------------
\subsection{Justificación. Condición necesaria:\cortetitulo los efectos desbordamiento}\label{sec:desbordamiento}

\begin{itemize}
\item \emph{Los 2 motivos para coordinar}\nocite{Bergin2008Palgrave}:
  \begin{itemize}
  \item los \concepto{efectos desbordamiento} (\emph{spillovers}) de la política económica, que son \concepto{externalidades}: quien decide no tiene en cuenta lo que su decisión hace a los demás\footnote{El libre comercio, la estabilidad financiera o un clima estable son \concepto{bienes públicos mundiales}: no rivales y difíciles de excluir. Cada país se beneficia de ellos los aporte o no, y tiende a aportar menos de lo que convendría a todos.} \vertema{3.A.23};
  \item el \concepto{reparto internacional del riesgo}: si los mercados no permiten a los hogares asegurarse contra las perturbaciones de su país, la política económica podría hacer ese papel (\ref{sec:corsetti-pesenti-extensiones}).
  \end{itemize}
\item \emph{Por qué los desbordamientos piden coordinación}:
  \begin{itemize}
  \item \emph{en la política fiscal}: ante una caída de la demanda mundial, cada país teme expandir solo, porque parte del estímulo se va en importaciones y, con tipo flexible, en apreciación. El empleo lo crean los vecinos y la deuda la paga quien expande. Cada uno espera a que lo haga otro, y la expansión se queda corta. Si todos se comprometen a expandir a la vez, lo que se escapa de cada uno vuelve de los demás;
  \item \emph{en la política monetaria}: una expansión deprecia la moneda y desvía la demanda hacia los bienes propios, a costa de los del vecino. Si todos lo hacen, nadie gana competitividad y todos acaban con más inflación: son las \concepto{depreciaciones competitivas}. Un acuerdo puede comprometer a todos a no hacerlo;
  \item \emph{su peso depende de la integración}: si las exportaciones son una parte pequeña del PIB, una depreciación apenas mueve la renta y los desbordamientos son pequeños. Y si los mercados no tratan igual la deuda de los distintos países, una expansión fiscal puede depreciar la moneda en vez de apreciarla, y el desbordamiento cambia de signo.
  \end{itemize}
\item \emph{Los canales de transmisión}:
  \begin{itemize}
  \item \emph{la renta}, por la cuenta corriente: parte de lo que gasta un país son importaciones, es decir, exportaciones de sus socios (el \concepto{efecto locomotora});
  \item \emph{el tipo de interés}, por la cuenta financiera: un país grande mueve el tipo de interés mundial y, con él, la inversión de los demás;
  \item \emph{el tipo de cambio}, por las 2 cuentas: la depreciación de uno es la apreciación del otro, que pierde competitividad;
  \item \emph{los precios relativos}: la relación real de intercambio, cuánto consume un país con lo que produce. Solo aparece en los modelos con microfundamentos (\ref{sec:corsetti-pesenti});
  \item \emph{los precios de los activos y el crédito}: el canal financiero, que la crisis de 2008 puso en primer plano.
  \end{itemize}
\item \emph{El canal financiero, en resumen}:
  \begin{itemize}
  \item los precios de los activos con riesgo, los flujos brutos de capital y el crédito se mueven a la vez en todo el mundo, y la política de la Reserva Federal da forma a ese \concepto{ciclo financiero global}: una subida de sus tipos baja los precios de los activos, el apalancamiento de los bancos y el crédito en todo el mundo\nocite{Rey2013Dilema,MirandaAgrippinoRey2020} \vertema{3.B.15};
  \item el canal no es el tipo de interés de cada país, sino los precios de los activos, el apalancamiento y la toma de riesgos; y un tipo de cambio flexible no aísla del todo: más que un trilema, un \emph{dilema} entre la libre movilidad del capital y una política monetaria propia \vertema{3.B.15};
  \item para la zona del euro, el BCE matiza que el motor principal es el apetito global por el riesgo y que no ha perdido el control de sus condiciones financieras\nocite{HabibVenditti2018};
  \item para la coordinación, es un argumento a favor: la Reserva Federal tiene un mandato nacional y no internaliza lo que su política hace fuera, salvo cuando vuelve sobre su propia economía.
  \end{itemize}
\item \emph{Demanda y oferta}: los desbordamientos de las políticas de oferta (reformas, comercio, fiscalidad, clima) son más claros y más duraderos que los de las políticas de demanda, y en ellos la coordinación es más frecuente y está más institucionalizada: la OMC, el acuerdo fiscal de la OCDE y el G-20, los acuerdos climáticos (\ref{sec:foros}).
\item Para ver los desbordamientos de las políticas de demanda, usamos 2 modelos de 2 países, uno con expectativas estáticas y otro con expectativas racionales.
\end{itemize}

\begin{anotaciones}
\begin{itemize}
\item \autor{Hélène Rey} explica el ciclo financiero global y la política monetaria en una entrevista en vídeo\nocite{Rey2013Video} y en su ponencia de Jackson Hole\nocite{Rey2013Dilema}.
\item La coordinación fue un gran tema de la macroeconomía internacional en los 80\nocite{BuiterMarston1985,BryantPortes1987}, y las cumbres del G-7 de Tokio (1986) y Venecia (1987) insistieron en la vigilancia multilateral. Pero, como apunta \textcite[nota~17]{Rey2013Dilema}, hasta hoy ni el análisis económico ni las declaraciones han tenido un efecto visible en la política monetaria de los bancos centrales.
\end{itemize}
\end{anotaciones}

% --------------------------------------------------------------------------
\subsubsection{Expectativas estáticas: \autor{Mundell}-\autor{Fleming} con 2 países}\label{sec:mundell-fleming-2}

\paragraph{a) Idea}
\begin{itemize}
\item El modelo de \textcite{Fleming1962Domestic} y \textcite{Mundell1963Capital} estudia un país pequeño: el tipo de interés mundial y la renta del resto del mundo le vienen dados \vertema{3.A.32}.
\item Para saber qué hace su política a los demás hay que dejar que el país sea grande y que la renta exterior responda. Es lo que hizo \textcite{Mundell1964Reply}, con 2 países grandes, A y B.
\item Es un modelo viejo y estático, pero muestra con sencillez los canales de la renta, el tipo de interés y el tipo de cambio, y por qué el signo del desbordamiento depende de la política y del régimen cambiario.
\end{itemize}

\paragraph{b) Modelo}
\subparagraph{i) Supuestos}
\begin{itemize}
\item Los del modelo de un país \vertema{3.A.32}:
  \begin{itemize}
  \item precios fijos ($P=P^*=1$), de modo que no hay expectativas de inflación y el tipo de cambio real se mueve con el nominal;
  \item expectativas estáticas sobre el tipo de cambio;
  \item activos nacionales y extranjeros sustitutos perfectos;
  \item movilidad perfecta del capital, $\beta_K\to\infty$, con $\beta_K$ la sensibilidad de los flujos de capital al diferencial de tipos (lleva subíndice porque $\beta$ es el factor de descuento);
  \item la condición de \autor{Marshall}-\autor{Lerner}: una depreciación real mejora el saldo comercial \vertema{3.B.12}.
  \end{itemize}
\item Y lo que distingue a esta variante:
  \begin{itemize}
  \item hay 2 países grandes, A (la zona del euro) y B (Estados Unidos), con tipo de cambio flexible entre ellos;
  \item la renta exterior de cada uno es la del otro y se determina dentro del modelo;
  \item con movilidad perfecta, los 2 tienen el mismo tipo de interés, que es el tipo mundial, $i=i^*$, y también se determina dentro del modelo.
  \end{itemize}
\end{itemize}

\subparagraph{ii) Desarrollo}
\begin{itemize}
\item \emph{Una expansión monetaria en A} (gráfico~\ref{fig:mundell-fleming-dos-paises}):
  \begin{itemize}
  \item en A, la LM se desplaza a la derecha y baja el tipo de interés; salen capitales hacia B, la moneda de A se deprecia y su IS se desplaza a la derecha. Como A es grande, el tipo de interés mundial también baja;
  \item en A, la renta sube, por el tipo y por la depreciación;
  \item en B actúan 2 fuerzas de signo contrario: la renta de A compra más importaciones a B, pero la moneda de B se aprecia y pierde competitividad. En el caso canónico domina la apreciación y la renta de B \emph{baja};
  \item es una política de \concepto{empobrecimiento del vecino} (\emph{beggar-thy-neighbour}): A gana renta a costa de B\nocite{Mundell1964Reply}. Siguiendo con el ejemplo, una expansión del BCE, por sí sola, subiría la renta de la zona del euro y bajaría la de Estados Unidos.
  \end{itemize}
\end{itemize}

\begin{figure}[H]
\centering
\caption{Una expansión monetaria en A con 2 países grandes, tipo de cambio flexible y movilidad perfecta}\label{fig:mundell-fleming-dos-paises}
\grafico[0.9\textwidth]{mundell-fleming-dos-paises}
\fuente{elaboración propia a partir de \textcite{Fleming1962Domestic} y \textcite{Mundell1963Capital}.}
\end{figure}

\begin{itemize}
\item \emph{Una expansión fiscal en A} (gráfico~\ref{fig:mundell-fleming-dos-paises-fiscal}): el gasto desplaza la IS de A a la derecha y sube su tipo de interés; entran capitales, la moneda de A se aprecia y la de B se deprecia y, como A es grande, el tipo de interés mundial sube.
  \begin{itemize}
  \item \emph{en A}, la renta sube, aunque menos que en una economía cerrada, porque la apreciación resta exportaciones netas;
  \item \emph{en B}, las 2 fuerzas van en el mismo sentido: A compra más importaciones a B y la depreciación da competitividad a B. La renta de B \emph{sube}: es el \concepto{efecto locomotora};
  \item la vuelta sobre A (B, más rico, compra más a A) la absorbe una apreciación adicional de la moneda de A, que mantiene $i=i^*$;
  \item siguiendo con el ejemplo, un estímulo fiscal en la zona del euro tira de la renta de Estados Unidos.
  \end{itemize}
\end{itemize}

\begin{figure}[H]
\centering
\caption{Una expansión fiscal en A con 2 países grandes, tipo de cambio flexible y movilidad perfecta}\label{fig:mundell-fleming-dos-paises-fiscal}
\grafico[0.9\textwidth]{mundell-fleming-dos-paises-fiscal}
\fuente{elaboración propia a partir de \textcite{Mundell1963Capital} y \textcite{Mundell1964Reply}.}
\end{figure}

\begin{anotaciones}
\textbf{El modelo de 2 países en 4 ecuaciones} (tipo flexible, movilidad perfecta, $P=P^*=1$; $q$ el tipo de cambio real de A, $m$ la propensión a importar, $\nu$ la sensibilidad del saldo comercial a $q$):
\[
\begin{aligned}
\text{IS}_A&:\ Y_A=\frac{\overline{A}_A-b\,i+m\,Y_B+\nu\,q}{1-c\,(1-t)+m}, &\qquad \text{LM}_A&:\ M_A=k\,Y_A-h\,i,\\
\text{IS}_B&:\ Y_B=\frac{\overline{A}_B-b\,i+m\,Y_A-\nu\,q}{1-c\,(1-t)+m}, &\qquad \text{LM}_B&:\ M_B=k\,Y_B-h\,i,
\end{aligned}
\]
con 4 incógnitas: $Y_A$, $Y_B$, $i$ y $q$. Las exportaciones de A son las importaciones de B.
\begin{itemize}
\item Sumando las 2 IS y las 2 LM sale un IS-LM del mundo, que fija la renta mundial, $Y_A+Y_B$, y el tipo de interés mundial, $i$.
\item Con el tipo mundial ya fijado, cada LM da la renta de su país: $Y_B=(M_B+h\,i)/k$.
  \begin{itemize}
  \item Si A expande el dinero, baja $i$ y, con $M_B$ igual, \emph{baja} $Y_B$: el empobrecimiento del vecino sale sin ambigüedad.
  \item Si A expande el gasto, sube $i$ y \emph{sube} $Y_B$: la locomotora. Si las 2 LM son iguales, la renta de B sube tanto como la de A, $\Delta Y_A=\Delta Y_B=(h/k)\,\Delta i$.
  \end{itemize}
\item Con tipo fijo, $q$ es dato y el dinero de cada país se ajusta con las reservas. Restando las 2 IS, la diferencia $Y_A-Y_B$ solo depende de $\overline{A}_A-\overline{A}_B$ y de $q$:
  \begin{itemize}
  \item la expansión monetaria baja $i$ y sube las 2 rentas por igual;
  \item la fiscal sube la renta mundial y el tipo $i$, y aumenta $Y_A-Y_B$: en B, el efecto depende de si pesa más la demanda de A o la subida del tipo de interés.
  \end{itemize}
\end{itemize}
\end{anotaciones}

\subparagraph{iii) Implicaciones}
\begin{itemize}
\item El signo del desbordamiento depende de la política (tabla~\ref{tab:mundell-fleming-2-completa}): con tipo flexible, la expansión fiscal es locomotora y la monetaria empobrece al vecino.
\item \emph{Para la coordinación}:
  \begin{itemize}
  \item la política monetaria con tipo flexible es un juego de suma negativa: cada uno quiere depreciar, y si todos lo hacen nadie gana;
  \item la fiscal es un juego de bienes públicos: cada uno quiere que tire el otro, y nadie tira;
  \item las 2 son dilemas del prisionero, y a ellos se vuelve en la deseabilidad (\ref{sec:deseabilidad}).
  \end{itemize}
\end{itemize}

\begin{table}[H]
\centering
\caption{Políticas expansivas en A con 2 países grandes y movilidad perfecta: efecto sobre la renta}\label{tab:mundell-fleming-2-completa}
\small
\begin{tabular}{>{\raggedright\arraybackslash}m{2.4cm}|>{\raggedright\arraybackslash}m{2.6cm}|>{\raggedright\arraybackslash}m{4.4cm}|>{\raggedright\arraybackslash}m{4.4cm}}
\rowcolor{tceeazulpalido}\textbf{Régimen} & \textbf{Política en A} & \textbf{Renta de A} & \textbf{Renta de B}\\
\hline
\multirow{2}{=}{Tipo flexible} & Monetaria & \textcolor{tceeverde}{Sube}; más cuanto menor es A & \textcolor{tceerojo}{Baja} en el caso canónico (domina la apreciación de B)\\
\hhline{~|-|-|-}
 & Fiscal & \textcolor{tceeverde}{Sube}; más cuanto mayor es A & \textcolor{tceeverde}{Sube} (efecto locomotora)\\
\hline
\multirow{2}{=}{Tipo fijo} & Monetaria & \textcolor{tceeverde}{Sube} & \textcolor{tceeverde}{Sube} (sube la oferta monetaria mundial)\\
\hhline{~|-|-|-}
 & Fiscal & \textcolor{tceeverde}{Sube} & \textcolor{tceedorado}{Ambiguo} (más demanda de A frente a un tipo de interés mundial más alto)\\
\hline
\end{tabular}
\fuente{elaboración propia a partir de \textcite{Mundell1963Capital} y \textcite{Mundell1964Reply}.}
\end{table}

\paragraph{c) Extensiones}
\begin{itemize}
\item \emph{El tamaño de A frente a B}:
  \begin{itemize}
  \item cuanto más pequeño es A, más se parece al país pequeño de \autor{Mundell}-\autor{Fleming}: su política monetaria gana eficacia, porque el tipo mundial apenas se mueve y toda la expansión actúa por la depreciación, y su política fiscal la pierde;
  \item cuanto más grande es A, más se parece a una economía cerrada: su política fiscal recupera eficacia, para A y para B, porque A arrastra el tipo mundial; la monetaria la pierde.
  \end{itemize}
\item \emph{Con tipo de cambio fijo} (y movilidad perfecta), el mundo es un solo mercado de dinero y la política de A se reparte entre los 2:
  \begin{itemize}
  \item una expansión monetaria en A baja el tipo de interés mundial; A pierde reservas, que gana B, y la renta \emph{sube en los 2 países};
  \item una expansión fiscal en A sube su renta y el tipo mundial; en B, el mayor gasto de A tira de sus exportaciones, pero el tipo de interés más alto frena su inversión: el efecto sobre B es \emph{ambiguo}.
  \end{itemize}
\item \emph{Con movilidad imperfecta}, las 2 fuerzas que actúan sobre B pierden potencia, pero las conclusiones no cambian en lo sustancial.
\end{itemize}

\paragraph{d) Evidencia empírica}
\begin{itemize}
\item La evidencia contradice el signo de la política monetaria:
  \begin{itemize}
  \item con tipos flexibles, una expansión de la Reserva Federal \emph{sube} la producción del resto del G-7, sobre todo porque baja el tipo de interés real mundial, y la balanza comercial pesa poco\nocite{Kim2001International};
  \item una subida de sus tipos \emph{baja} la producción en casi todos los países, también en los que flotan\nocite{DedolaRivoltaStracca2017};
  \item tras 2008, la expansión de los países avanzados empujó capital hacia los emergentes, y el movimiento se invirtió cuando la Reserva Federal empezó a subir tipos, en diciembre de 2015\nocite{Fed2015FOMCdic}, y a reducir su balance, en octubre de 2017\nocite{Fed2017Balance}.
  \end{itemize}
\item Para explicarlo hacen falta canales que \autor{Mundell}-\autor{Fleming} no tiene: los precios relativos, que añade el modelo siguiente, y el canal financiero.
\end{itemize}

\paragraph{e) Valoración}
\begin{itemize}
\item \emph{Ventajas}: muestra con sencillez los canales de la interdependencia y por qué el signo del desbordamiento depende de la política y del régimen.
\item \emph{Inconvenientes}:
  \begin{itemize}
  \item es estático, con precios fijos y sin microfundamentos: no dice nada del bienestar, solo de la renta, y está expuesto a la crítica de \textcite{Lucas1976Econometric}\footnote{\autor{Robert Lucas} recibió el Premio Nobel de Economía en 1995 «por haber desarrollado y aplicado la hipótesis de las expectativas racionales, transformando así el análisis macroeconómico y profundizando en la comprensión de la política económica»\nocite{NobelPrize1995Lucas}.} \vertema{3.A.6};
  \item las políticas son exógenas: no tiene en cuenta cómo responde cada país a la política del otro, que es justo lo que estudia la teoría de juegos (\ref{sec:deseabilidad});
  \item no dice si coordinar compensa: solo que hay desbordamientos.
  \end{itemize}
\item Las 2 primeras críticas las responde la nueva macroeconomía abierta, con \autor{Corsetti} y \autor{Pesenti}.
\end{itemize}

% --------------------------------------------------------------------------
\subsubsection{Expectativas racionales:\cortetitulo la nueva macroeconomía abierta de \autor{Corsetti} y \autor{Pesenti}}\label{sec:corsetti-pesenti}

\paragraph{a) Idea}
\begin{itemize}
\item La \concepto{nueva macroeconomía abierta} (\emph{new open economy macroeconomics}, NOEM) nace con el modelo \emph{Redux} de \textcite{ObstfeldRogoff1995Redux}: modelos de equilibrio general dinámicos con agentes que optimizan, competencia monopolística y rigideces nominales, en 2 países \vertema{3.A.7} \vertema{3.B.14}. Son la versión abierta del modelo neokeynesiano con microfundamentos que sistematizó \autor{Michael Woodford}.\footnote{\autor{Michael Woodford} recibió el Premio Nobel de Economía en 2026, con \autor{Robert E. Hall}, «por los fundamentos microeconómicos de la estabilización macroeconómica»\nocite{Nobel2026Prensa}: es el marco en el que se preguntan \autor{Corsetti} y \autor{Pesenti} si conviene coordinar.}
\item Son los modelos de trabajo de bancos centrales y organismos, como el SIGMA de la Reserva Federal\nocite{ErcegGuerrieriGust2005}, el GEM del FMI\nocite{Bayoumi2004GEM} o los de la Comisión Europea.\footnote{La Comisión Europea usa QUEST, un modelo de esta familia para la zona del euro y sus Estados miembros\nocite{CE2026QUEST}, y el modelo global multipaís GM\nocite{CE2026GM}.} Su versión para una economía pequeña es la de \textcite{GaliMonacelli2005Monetary} \vertema{3.A.32}; aquí interesan los de 2 países.
\item \textcite{CorsettiPesenti2009} reducen uno de esos modelos a 3 curvas en el plano del empleo y el consumo, para enseñarlo como el IS-LM. Aporta 2 cosas que \autor{Mundell}-\autor{Fleming} no tiene:
  \begin{itemize}
  \item un \emph{canal nuevo}, la relación real de intercambio;
  \item un \emph{criterio de bienestar}: la utilidad de los hogares, con la que se puede preguntar si coordinar mejora a los 2 países. Por eso son modelos de coordinación con microfundamentos, a diferencia del diagrama de \autor{Hamada} (\ref{sec:hamada}), que parte de funciones de pérdida postuladas.
  \end{itemize}
\end{itemize}

\paragraph{b) Modelo}
\subparagraph{i) Supuestos}
\begin{itemize}
\item 2 países, A y B (las variables de B llevan subíndice $B$, salvo su tono monetario, $\mu^*$, que entra en el tipo de cambio), con un hogar representativo cuya utilidad crece con el consumo y decrece con el trabajo, $U=\ln C-\chi\,\ell$, con $\chi$ el peso de la desutilidad del trabajo y $\ell$ el empleo.
\item En el plano $(\ell,C)$, las curvas de indiferencia son crecientes y el bienestar sube hacia arriba y a la izquierda: más consumo con menos trabajo.
\item Competencia monopolística entre variedades, con elasticidad de sustitución $\sigma>1$\nocite{DixitStiglitz1977Monopolistic} \vertema{3.A.18}; tecnología lineal en el trabajo, con productividad $Z$.
\item Sin inversión ni gasto público: lo que se produce se consume. El consumo es una cesta con la mitad de bienes de cada país.
\item Precios fijados por adelantado un periodo (rigidez nominal) y mercados financieros completos.
\item Un banco central benevolente, que maximiza la utilidad del hogar de su país, con un instrumento, el tono de la política monetaria, $\mu$: el gasto nominal que sostiene (aquí $\mu$ no es el margen).
\end{itemize}

\subparagraph{ii) Desarrollo}
\begin{itemize}
\item El modelo se resume en 3 curvas en el plano $(\ell,C)$ y una ecuación para el tipo de cambio:
\funcionclave[eq:corsetti-pesenti]{\text{AD: }C=\frac{\mu}{P},\qquad \text{AS: }C=Z\,\ell\,\tau,\qquad \text{NR: }\overline{\ell}=\frac{\sigma-1}{\sigma\,\chi},\qquad E=\frac{\mu}{\mu^*}}
  \begin{itemize}
  \item \emph{AD}, la demanda agregada: el gasto nominal $\mu$ dividido por el nivel de precios. Es horizontal en el plano $(\ell,C)$: con precios fijos, el consumo lo fija la política monetaria, que es el rasgo keynesiano del modelo;
  \item \emph{AS}, la tecnología: un rayo desde el origen, cuya pendiente es la productividad $Z$ por $\tau$, el término de la \concepto{relación real de intercambio}: cuánto consumo compra una unidad de lo producido. Si los bienes de A se abaratan frente a los de B, $\tau$ baja y la AS gira hacia abajo: para consumir lo mismo hay que trabajar más\footnote{Formalmente, $\tau\equiv\big[\tfrac{P}{2}\big(\tfrac{1}{P_H}+\tfrac{1}{E\,P^*_H}\big)\big]^{-1}$, con $P_H$ el precio en A de lo producido en A y $P^*_H$ su precio en B, en moneda de B: $1/\tau$ son las unidades de producción de A que hacen falta para comprar una cesta de consumo de A. En la economía cerrada, $\tau=1$\nocite{CorsettiPesenti2009}. Lo que los autores llaman AS no es la oferta agregada de los manuales: es la restricción tecnológica.};
  \item \emph{NR}, la tasa natural: el empleo con precios flexibles, $\overline{\ell}$, una vertical. No depende de la productividad y sí del poder de mercado: cuanto menor es $\sigma$ (variedades más diferenciadas, mayor margen), menos se trabaja; con competencia perfecta ($\sigma\to\infty$) sería el empleo eficiente, $1/\chi$;
  \item \emph{el tipo de cambio} es el cociente de los tonos monetarios: con mercados completos, una política más expansiva en A deprecia su moneda \vertema{3.B.14}.
  \end{itemize}
\item \emph{El equilibrio}, con precios rígidos, está donde se cortan la AD y la AS; la NR marca el empleo al que volvería la economía con precios flexibles, el que un banco central benevolente quiere recuperar ante una perturbación.
\item Por el poder de mercado, en el equilibrio inicial se trabaja menos de lo eficiente: trabajar algo más y consumir lo que se produce subiría el bienestar. Es la distorsión monopolística, que el modelo hereda de \autor{Dixit} y \autor{Stiglitz}.\footnote{\autor{Joseph Stiglitz} recibió el Premio Nobel de Economía en 2001, con \autor{George Akerlof} y \autor{Michael Spence}, «por sus análisis de los mercados con información asimétrica»\nocite{Nobel2001AkerlofSpenceStiglitz}.}
\item \emph{Una subida de la productividad en A} ($Z\uparrow$), que gira hacia arriba la AS de A (gráfico~\ref{fig:corsetti-pesenti}, panel a; A se comporta, a corto plazo, como una economía cerrada):
  \begin{itemize}
  \item \emph{con precios rígidos}, ni los precios ni el tipo de cambio se mueven, y la AD de A no cambia. A pasa de $O$ a $R$: consume lo mismo con menos trabajo, por debajo del empleo natural (una brecha de producción negativa), y B no se entera;
  \item \emph{con precios flexibles}, los bienes de A se abaratan y la perturbación se reparte entre los 2 países por la relación real de intercambio: A llega a $F$, sobre su NR, con más consumo, y B también consume más con el mismo empleo, porque se abarata lo que importa;
  \item la brecha invita a responder con la política monetaria: el banco central de A sube $\mu$ (AD$'$) y lleva a A de $R$ a $F$, el equilibrio con precios flexibles. Es esa respuesta la que lleva la perturbación a B: el desbordamiento depende de la política, no de la perturbación.
  \end{itemize}
\item \emph{Lo que esa expansión hace a B} depende de en qué moneda fijan sus precios de exportación las empresas, porque eso decide cuánto del movimiento del tipo de cambio pasa a los precios de importación (el \concepto{traspaso}, \emph{pass-through}) \vertema{3.A.37}:
  \begin{itemize}
  \item en la moneda del productor (\emph{producer currency pricing}, PCP): traspaso completo;
  \item en la moneda del comprador (\emph{local currency pricing}, LCP): traspaso nulo a corto plazo;
  \item todos en la moneda de uno de los 2 países, el «dólar» (\emph{dollar pricing}): traspaso asimétrico.
  \end{itemize}
\item \emph{Con PCP} (panel b):
  \begin{itemize}
  \item la expansión deprecia la moneda de A ($E=\mu/\mu^*$ sube) y la depreciación pasa entera a los precios de importación: los bienes de A se abaratan en B y los de B se encarecen en A, como en \autor{Mundell}-\autor{Fleming}. La relación de intercambio de A empeora ($\tau\downarrow$) y la de B mejora ($\tau_B\uparrow$);
  \item en B, la AS gira hacia arriba y las importaciones más baratas bajan su índice de precios y suben su AD: B pasa de $O$ a $H$, sobre su NR, con más consumo y el mismo empleo. El desbordamiento es \emph{positivo};
  \item el tipo de cambio hace lo que harían los precios si fueran flexibles: la política que cierra la brecha en A reproduce el equilibrio con precios flexibles en los 2 países. Por la misma razón, el tipo de cambio salta de una vez a su nuevo valor, sin sobrerreacción \vertema{3.B.14};
  \item frente a \autor{Mundell}-\autor{Fleming}, el signo se invierte: allí la expansión monetaria empobrece al vecino porque le quita demanda; aquí le beneficia porque le abarata lo que compra. El canal es la relación real de intercambio, no las exportaciones netas.
  \end{itemize}
\item \emph{Con LCP} (panel b): cada empresa fija en la moneda del comprador el precio de lo que vende en cada país, de modo que a corto plazo todos los precios de consumo del mundo están fijos, se mueva como se mueva el tipo de cambio.
  \begin{itemize}
  \item La depreciación de A no abarata los bienes de A en B ni los encarece en A: lo que cambia son los márgenes. El exportador de A cobra en B el mismo precio en moneda de B, que ahora vale más en moneda de A: la relación de intercambio de A \emph{mejora} y la de B empeora ($\tau_B\downarrow$), justo al revés que con PCP.
  \item En A, la AD sube tanto como $\mu$, porque no hay inflación importada, y A consume más que con PCP.
  \item En B, la AD no se mueve (sus precios de consumo están fijos en su moneda) y la AS gira hacia abajo: B pasa de $O$ a $L$, con el mismo consumo y más trabajo, fuera de su empleo natural. Su bienestar baja: la expansión de A \emph{empobrece al vecino}.
  \end{itemize}
\end{itemize}

\begin{figure}[H]
\centering
\caption{Perturbación de productividad y respuesta monetaria en el modelo de \autor{Corsetti} y \autor{Pesenti}}\label{fig:corsetti-pesenti}
\grafico[0.9\textwidth]{corsetti-pesenti}
\fuente{elaboración propia a partir de \textcite{CorsettiPesenti2009}. $\mu$ es aquí el tono de la política monetaria, no el margen.}
\end{figure}

\begin{itemize}
\item \emph{Por qué importa el caso LCP}:
  \begin{itemize}
  \item añade una rigidez nominal más, la de los precios de importación, que se suma a la de los precios internos;
  \item tiene respaldo empírico: a corto plazo, los precios de importación apenas responden al tipo de cambio, sobre todo cuando el comercio se factura en dólares, que es lo más frecuente\nocite{Gopinath2020DCP} \vertema{3.A.37}.
  \end{itemize}
\item \emph{Sus consecuencias para los desbordamientos}:
  \begin{itemize}
  \item un movimiento del tipo de cambio del mismo tamaño que con PCP apenas cambia el comercio a corto plazo, porque no hay \concepto{desvío del gasto} (\emph{expenditure switching}) entre países: lo que se desvía, dicen los autores, es la carga de trabajo necesaria para sostener el consumo mundial;
  \item como los precios de consumo no absorben la expansión monetaria, el tipo de cambio tiene que moverse más que con PCP para equilibrar el mercado de dinero: con fijación de precios por mercados (\emph{pricing to market}) vuelve la \concepto{sobrerreacción} del tipo de cambio de \autor{Dornbusch}\nocite{BettsDevereux2000} \vertema{3.B.14};
  \item la política monetaria pasa de locomotora a empobrecimiento del vecino: los desbordamientos dependen de cómo se fijan los precios.
  \end{itemize}
\item \emph{Con facturación en dólares}, todo el comercio se factura en la moneda de uno de los 2 países, el «dólar» (en el ejemplo, la de Estados Unidos, B), no en una tercera moneda. Cuando el emisor cambia su política, sus precios de importación no se mueven y queda aislado de lo que pasa fuera, mientras el otro país recibe todo su efecto: es la versión de 2 países del \concepto{paradigma de la moneda dominante}\nocite{Gopinath2020DCP}. El detalle gráfico está en el anexo~\ref{anexo:corsetti-pesenti}.
\end{itemize}

\subparagraph{iii) Implicaciones}
\begin{itemize}
\item \emph{¿Hay ganancias de coordinar?} Los autores comparan la política óptima de cada banco central por su cuenta (el equilibrio de \autor{Nash}) con la que elegirían juntos, maximizando la media de los 2 bienestares\nocite{CorsettiPesenti2009}:
  \begin{itemize}
  \item \emph{con PCP}, ninguna: la política óptima mira solo hacia dentro (cerrar la brecha propia) y el tipo de cambio hace el resto. Es la versión extrema del argumento de \textcite{Friedman1953Case}: aun con precios rígidos, el tipo de cambio da los precios relativos correctos;
  \item \emph{con LCP}, tampoco, aunque haya desbordamientos: cada banco central estabiliza una media de los costes de las empresas de los 2 países, para no hacer variar los márgenes de los exportadores, y los 2 eligen la misma política, de modo que el tipo de cambio no se mueve. Lo que se ganaría coordinando, un tipo de cambio estable, ya se consigue sin coordinar\nocite{CorsettiPesenti2009,DevereuxEngel2003};
  \item el efecto de empobrecer al vecino existe, pero ningún banco central lo busca: en el óptimo los 2 eligen la misma política y el tipo de cambio no se mueve\nocite{DevereuxEngel2003};
  \item \emph{con facturación en dólares}, sí: el bienestar mundial sube si se coordina, pero solo tiene que cambiar su política el país que emite el dólar. La contribución a la cooperación es unilateral, y cabe preguntarse qué incentivo tiene ese país a firmar un acuerdo.
  \end{itemize}
\item En los 2 primeros casos, «poner orden en la propia casa» es la mejor regla. El tercero, que es el que más se parece a un mundo que factura en dólares, deja abierta la puerta.
\item Es la distinción del título del apartado: con LCP hay efectos desbordamiento (B queda en $L$, fuera de su empleo natural), pero no basta para que coordinar mejore el resultado. La condición suficiente es otra (\ref{sec:deseabilidad}).
\end{itemize}

\paragraph{c) Extensiones}\label{sec:corsetti-pesenti-extensiones}
\begin{itemize}
\item \emph{Una reforma estructural en A} que reduce el poder de mercado de sus empresas sube $\sigma$ (variedades más sustituibles, menor margen) y, con él, el empleo natural de A:
  \begin{itemize}
  \item A produce más, sus bienes se abaratan y su relación de intercambio empeora ($\tau\downarrow$), mientras la de B mejora ($\tau_B\uparrow$);
  \item A trabaja y consume más, aunque el empeoramiento de la relación de intercambio le quita parte de la ganancia; B consume más con el mismo trabajo;
  \item es un desbordamiento \emph{positivo} de una política de \emph{oferta}: las reformas que reducen las distorsiones de un país benefician también a sus socios\nocite{CorsettiPesenti2009}. El gráfico está en el anexo~\ref{anexo:corsetti-pesenti}.
  \end{itemize}
\item \emph{El reparto del riesgo}: si una caída de la productividad golpea solo a A y los mercados no permiten asegurarse, la política podría hacer de seguro. Una contracción monetaria en A, por ejemplo, aprecia su moneda, encarece sus exportaciones y le transfiere renta desde B. En manos de un país solo, es la tentación de empobrecer al vecino; en manos de un coordinador, una forma de repartir el golpe\nocite{Bergin2008Palgrave} \vertema{3.B.12}.
  \begin{itemize}
  \item El problema es que una política con un solo instrumento no puede corregir a la vez 2 distorsiones, la rigidez de precios y la falta de seguro: lo que pide una no tiene por qué ser lo que pide la otra\nocite{ObstfeldRogoff2002Global}.
  \item La integración reduce este motivo: con más comercio, la relación de intercambio ya reparte el golpe (los bienes de A, más escasos, se encarecen), y con más mercados financieros, los hogares se aseguran solos. Paradójicamente, la globalización puede pedir \emph{menos} coordinación, no más\nocite{Bergin2002FRBSF}.
  \end{itemize}
\item \emph{Los bienes intermedios importados y las cadenas globales de valor}: si lo que se importa entra en la producción propia, el tipo de cambio afecta a los costes y los desbordamientos son mayores, sobre todo en las economías pequeñas que ensamblan para exportar \vertema{3.B.18}.
\item \emph{Los bienes no comerciables}: si las perturbaciones golpean de forma distinta a los sectores de bienes comerciables y no comerciables, el tipo de cambio no basta para ajustar los precios relativos, y la política óptima ya no es mirar solo hacia dentro\nocite{CorsettiPesenti2009}.
\end{itemize}

\paragraph{d) Evidencia empírica}
\begin{itemize}
\item \emph{Cómo se fijan los precios} es una cuestión empírica, y la respuesta ha cambiado: el traspaso a los precios de importación es bajo a corto plazo y la mayor parte del comercio se factura en dólares, no en la moneda del exportador ni en la del comprador\nocite{Gopinath2020DCP}. El mundo se parece más al caso del dólar que a los otros 2.
\item \emph{Cuánto se gana coordinando}: en los modelos de la NOEM, poco\nocite{ObstfeldRogoff2002Global}; con bienes intermedios importados y no comerciables, la ganancia puede llegar a ser apreciable para algunos países, hasta el 2,2~\% del consumo en algún caso\nocite{Tchakarov2004}.
\end{itemize}

\paragraph{e) Valoración}
\begin{itemize}
\item \emph{Ventajas}: microfundamentos, un criterio de bienestar y un canal nuevo, la relación de intercambio, en un gráfico de 3 curvas.
\item \emph{Limitaciones}: los resultados dependen de supuestos difíciles de contrastar (en qué moneda se fijan los precios, qué activos hay, preferencias logarítmicas que hacen coincidir \autor{Nash} y cooperación); con preferencias más generales, la política de cada uno y la cooperativa ya no coinciden\nocite{BenignoBenigno2006}.
\item Y no explica lo que se ve en la práctica: los gobiernos coordinan poco aunque haya desbordamientos. Para eso hay que pasar de los efectos a las decisiones.
\end{itemize}

\begin{anotaciones}
El artículo completo, con todos los casos y sus gráficos, está en la versión de trabajo del NBER\nocite{CorsettiPesenti2005}; la versión publicada es la de \textcite{CorsettiPesenti2009}.
\end{anotaciones}

\begin{itemize}
\item Recapitulando la primera condición: la política de un país grande desborda sus fronteras, pero el signo depende del modelo y de cómo se fijan los precios, y que haya desbordamientos no garantiza que coordinar mejore el resultado. Esto nos lleva a la segunda condición.
\end{itemize}

% --------------------------------------------------------------------------
\subsection{Deseabilidad. Condición suficiente:\cortetitulo las ganancias. La teoría de juegos}\label{sec:deseabilidad}

\begin{itemize}
\item Una vez vistos los desbordamientos, la condición necesaria, pasamos a la suficiente: que coordinar dé a todos un resultado mejor que actuar cada uno por su cuenta.
\item Los modelos de coordinación usan la \concepto{teoría de juegos}, porque estudian decisiones con \concepto{interdependencia estratégica}: lo que gana cada gobierno depende de lo que haga el otro, y los 2 lo saben \vertema{3.A.14}.
\item Casi todos son variantes del \concepto{dilema del prisionero}: si los 2 países cooperan, están mejor que si no lo hacen, pero no cooperar es la mejor respuesta de cada uno a lo que haga el otro.
\end{itemize}

% --------------------------------------------------------------------------
\subsubsection{El diagrama de \autor{Hamada}}\label{sec:hamada}

\paragraph{a) Idea}
\begin{itemize}
\item \textcite{Hamada1976Strategic} llevó la teoría de juegos a la política monetaria entre 2 países con un diagrama que todavía se usa\footnote{Es \autor{Koichi Hamada}, que lo presentó ya en 1974\nocite{Hamada1974Alternative}. No hay que confundirlo con \autor{Robert Hamada}, el de la beta apalancada de las finanzas (1972).}; \textcite{Corden1994Economic} lo expone con el tipo de cambio como objetivo, que es la versión que contamos.
\item Permite ver en un solo gráfico por qué el equilibrio sin coordinar es ineficiente y por qué la solución cooperativa es frágil.
\item Es el problema de «¿quién le pone el cascabel al gato?»: los 2 países saben lo que les conviene a ambos, pero cada uno espera a que se mueva el otro, porque el que se mueve solo carga con el coste.
\end{itemize}

\paragraph{b) Modelo}
\subparagraph{i) Supuestos}
\begin{itemize}
\item 2 países, A y B, cada uno con un instrumento, su oferta monetaria, $M_A$ y $M_B$.
\item Cada país tiene una disyuntiva entre inflación, $\pi_i$, y paro, $U_i$ (la curva de \autor{Phillips}) \vertema{3.A.28}, y una pérdida con la forma de la de \autor{Barro} y \autor{Gordon} \vertema{3.A.36}; B es más contrario a la inflación que A ($\gamma_B>\gamma_A$).
\item El tipo de cambio lo fija el dinero relativo, como en el modelo monetario \vertema{3.B.14}: $E=f(M_A/M_B)$, con $f'>0$. Una expansión relativa de A deprecia su moneda, lo que le sube la inflación (las importaciones son más caras) y le baja el paro (exporta más).
\item Se cumple la paridad no cubierta de intereses: los tipos de interés de los 2 países difieren en la depreciación esperada \vertema{3.B.13}.
\item Los 2 países comparten un objetivo de tipo de cambio, $E_0$: el tipo de cambio compatible con los objetivos internos de cada uno es el mismo.
\item El juego es estático y se juega una vez.
\end{itemize}

\subparagraph{ii) Desarrollo}
\begin{itemize}
\item \emph{La pérdida de cada país}, $i\in\{A,B\}$, crece con la distancia del paro y de la inflación a sus objetivos, $U^*_i$ y $\pi^*_i$, con un peso $\gamma_i$ para la inflación:
\funcionclave[eq:hamada]{L_i=(U_i-U^*_i)^2+\gamma_i\,(\pi_i-\pi^*_i)^2,\qquad i\in\{A,B\}}
\item \emph{El problema de cada país}, por ejemplo A, toma como dada la política de B:
$$\optim{\min}{M_A}{L_A=(U_A-U^*_A)^2+\gamma_A(\pi_A-\pi^*_A)^2}{\pi_A=\pi_A(\underset{+}{E},\underset{+}{M_A})\\ U_A=U_A(\underset{-}{E},\underset{-}{M_A})\\ E=f(M_A/M_B)}$$
  \begin{itemize}
  \item la primera restricción dice que la depreciación y la expansión monetaria suben la inflación; la segunda, que bajan el paro;
  \item su solución, para cada $M_B$, es la \concepto{función de reacción} de A, $M_A=R_A(M_B)$; la de B es $M_B=R_B(M_A)$.
  \end{itemize}
\item \emph{El gráfico} (gráfico~\ref{fig:hamada}), con $M_A$ en el eje horizontal y $M_B$ en el vertical:
  \begin{itemize}
  \item cada rayo desde el origen es un tipo de cambio: cuanto más inclinado, mayor $M_B$ frente a $M_A$ y más apreciada la moneda de A;
  \item cada país tiene un \concepto{punto de felicidad}, el de pérdida nula, en el que alcanza a la vez sus objetivos de paro y de inflación: $F_A$ para A, más lejos del origen porque A prefiere una política más laxa en los 2 países, y $F_B$ para B. Los 2 están en el rayo $E_0$;
  \item las curvas de indiferencia de cada país son círculos alrededor de su punto de felicidad, peores cuanto más lejos;
  \item con curvas circulares, la función de reacción de A es la vertical que pasa por $F_A$: haga lo que haga B, a A le conviene el mismo $M_A$, el que le deja en el círculo más cercano a $F_A$. La de B es la horizontal que pasa por $F_B$.
  \end{itemize}
\item \emph{El equilibrio no cooperativo} es el corte de las 2 funciones de reacción, el punto $N$: un \concepto{equilibrio de Nash}, porque cada país hace lo mejor que puede dado lo que hace el otro \vertema{3.A.14}.
  \begin{itemize}
  \item No es un \concepto{óptimo de Pareto}: por $N$ pasan una curva de indiferencia de cada país que forman una lente, y en cualquier punto de la lente los 2 están mejor.
  \item En $N$, el tipo de cambio $E_N$ deja la moneda de A demasiado depreciada para los 2.
  \end{itemize}
\item \emph{La solución cooperativa}: los 2 países eligen juntos sus políticas para minimizar una media de sus pérdidas, $\omega\,L_A+(1-\omega)\,L_B$, con un peso $\omega\in(0,1)$ que refleja el poder negociador de A.
  \begin{itemize}
  \item Sus soluciones, al variar $\omega$, son los puntos de tangencia de las curvas de indiferencia: la \concepto{curva de contrato}, que con círculos es el segmento que une $F_A$ y $F_B$.
  \item La parte de la curva de contrato que queda dentro de la lente es el \concepto{núcleo}: los acuerdos que mejoran a los 2 frente a $N$. El punto concreto depende del poder negociador.
  \item Como los 2 puntos de felicidad están en el mismo rayo, todo el núcleo está sobre $E_0$: el tipo de cambio cooperativo es el objetivo común.
  \end{itemize}
\end{itemize}

\begin{figure}[H]
\centering
\caption{El diagrama de \autor{Hamada}: el equilibrio de \autor{Nash} frente a la solución cooperativa}\label{fig:hamada}
\grafico[0.7\textwidth]{hamada}
\fuente{elaboración propia a partir de \textcite{Hamada1976Strategic} y \textcite{Corden1994Economic}.}
\end{figure}

\subparagraph{iii) Implicaciones}
\begin{itemize}
\item Con desbordamientos, actuar cada uno por su cuenta es ineficiente, y coordinar mejora a los 2: es el argumento teórico de la coordinación.
\item Pero la solución cooperativa \emph{no es un equilibrio de Nash}: si A se compromete con su parte del acuerdo, a B le conviene volver a su función de reacción. Cada país tiene incentivos a comportarse como un \concepto{polizón} (\emph{free rider}).
\item Es un dilema del prisionero en versión continua: con curvas circulares, cada país tiene incluso una estrategia dominante, como en el dilema.
\item Que coordinar sea mejor no basta: hace falta que la cooperación se sostenga.
\end{itemize}

\paragraph{c) Extensiones}
\begin{itemize}
\item \emph{Si uno de los 2 decide primero}, a la \autor{Stackelberg} \vertema{3.A.19}, el líder elige el punto de la función de reacción del otro que más le conviene. Con curvas circulares, el punto de la horizontal de B más cercano a $F_A$ es $N$: el liderazgo no cambia nada. Con curvas elípticas, las funciones de reacción tienen pendiente, el liderazgo sí cambia el resultado y el núcleo sigue siendo la parte de la curva de contrato dentro de la lente\nocite{Hamada1976Strategic}.
\item \emph{Si el juego se repite}, la cooperación puede sostenerse sin contrato (\ref{sec:repeticion}).
\end{itemize}

\paragraph{d) Evidencia empírica}
\begin{itemize}
\item \textcite{OudizSachs1984} cuantificaron por primera vez un juego de este tipo, con modelos econométricos keynesianos de Estados Unidos, Alemania y Japón. La ganancia de coordinar fue modesta: para Estados Unidos, el equivalente a medio punto del PNB al año, algo parecido para Alemania y algo más para Japón.
\item La razón es el poco comercio: en 1982 las exportaciones de mercancías de Estados Unidos a la Comunidad Europea eran solo el 1,6~\% de su PNB\nocite{OudizSachs1984}.
\item Los repasos posteriores de la evidencia llegaron a la misma conclusión: las ganancias de coordinar son pequeñas\nocite{McKibbin1997Empirical,OstryGhosh2013}.
\end{itemize}

\paragraph{e) Valoración}
\begin{itemize}
\item \emph{Ventajas}: muestra en un gráfico la diferencia entre el equilibrio de \autor{Nash} y el óptimo de \autor{Pareto}, y por qué la solución cooperativa es frágil.
\item \emph{Limitaciones}:
  \begin{itemize}
  \item las curvas circulares y el objetivo común de tipo de cambio son supuestos muy simplificadores;
  \item coordinar solo a los bancos centrales puede ser contraproducente si no tienen credibilidad frente al sector privado\nocite{Rogoff1985Counterproductive}; el argumento se desarrolla con los obstáculos (\ref{sec:obstaculos});
  \item las ganancias medidas son pequeñas: \textcite{BlanchardOstryGhosh2013} llamaron a la coordinación «el monstruo del lago Ness», de la que se habla mucho más de lo que se la ve.
  \end{itemize}
\end{itemize}

% --------------------------------------------------------------------------
\subsubsection{La repetición: cómo se sostiene la cooperación}\label{sec:repeticion}

\paragraph{a) Idea}
\begin{itemize}
\item En el diagrama de \autor{Hamada}, la solución cooperativa no se sostiene si el juego se juega una vez. Pero los gobiernos no se ven una sola vez: coinciden en todas las cumbres.
\item Si el juego se repite, la amenaza de dejar de cooperar en el futuro puede sostener hoy la cooperación sin necesidad de un contrato.
\end{itemize}

\paragraph{b) Modelo}
\subparagraph{i) Supuestos}
\begin{itemize}
\item \emph{El juego de la locomotora}\nocite{Frankel2016Coordination}: ante una recesión mundial, la zona del euro (A) y Estados Unidos (B) deciden si expandir el gasto o no, con los mismos pagos que la guerra de precios del dilema del prisionero \vertema{3.A.14} (tabla~\ref{tab:locomotora}):
  \begin{itemize}
  \item si los 2 expanden, la economía mundial sale de la recesión y nadie empeora su balanza comercial: 5 cada uno;
  \item si expande uno solo, paga la locomotora de otro: el que expande se queda con el déficit y la deuda (0) y el otro se lleva el tirón sin pagarlo (6);
  \item si no expande ninguno, la recesión sigue: 1 cada uno.
  \end{itemize}
\item El juego se repite sin un final conocido, y cada gobierno descuenta los pagos futuros con un factor $\beta\in(0,1)$.
\item Cada gobierno sigue la \concepto{estrategia del gatillo} de \autor{James W. Friedman} (1971)\nocite{Friedman1971Supergames}: coopera mientras el otro coopere y, si el otro se desvía una vez, deja de cooperar para siempre.
\end{itemize}

\begin{table}[H]
\centering
\caption{El juego de la locomotora: la zona del euro (filas) y Estados Unidos (columnas) ante una recesión mundial}\label{tab:locomotora}
\begin{tabular}{>{\centering\arraybackslash}m{2.6cm}|>{\centering\arraybackslash}m{2.6cm}|>{\centering\arraybackslash}m{2.6cm}|}
\multicolumn{1}{c}{} & \multicolumn{1}{c}{\textbf{Expandir}} & \multicolumn{1}{c}{\textbf{No expandir}}\\
\hhline{~|-|-|}
\textbf{Expandir} & $(5,5)$ & $(0,\underline{6})$\\
\hhline{~|-|-|}
\textbf{No expandir} & $(\underline{6},0)$ & \cellcolor{tceesalvia}$(\underline{1},\underline{1})$\\
\hhline{~|-|-|}
\end{tabular}
\fuente{elaboración propia a partir de \textcite{Frankel2016Coordination}, con los pagos del dilema del prisionero de la guerra de precios. Subrayadas, las mejores respuestas; sombreado, el equilibrio de \autor{Nash}.}
\end{table}

\subparagraph{ii) Desarrollo}
\begin{itemize}
\item \emph{Si el juego se jugara una vez}, no expandir sería una estrategia dominante y (No expandir, No expandir) el único equilibrio de \autor{Nash}, con 1 cada uno.
\item \emph{En el juego repetido}, cada gobierno compara lo que gana desviándose hoy con lo que pierde si el otro deja de cooperar para siempre. Con $u^{\mathrm C}$ el pago por periodo de cooperar, $u^{\mathrm D}$ el de desviarse y $u^{\mathrm N}$ el del equilibrio de \autor{Nash} que sirve de castigo, cooperar es un equilibrio si los gobiernos son lo bastante pacientes:
\funcionclave[eq:gatillo]{\beta\;\geq\;\frac{u^{\mathrm{D}}-u^{\mathrm{C}}}{u^{\mathrm{D}}-u^{\mathrm{N}}}}
  \begin{itemize}
  \item el numerador es lo que se gana hoy desviándose; el denominador, lo que se pierde cada periodo después;
  \item volvamos a la locomotora: desviarse gana $6-5=1$ hoy y pierde $5-1=4$ cada periodo después; cooperar compensa si $\beta\geq1/5$.
  \end{itemize}
\end{itemize}

\subparagraph{iii) Implicaciones}
\begin{itemize}
\item La cooperación puede sostenerse sin contrato: es la lógica del teorema de tradición oral (\emph{folk theorem}) \vertema{3.A.14} y de la colusión tácita entre empresas \vertema{3.A.19}.
\item \emph{Qué es aquí el factor de descuento}: la paciencia de los gobiernos, que depende del horizonte electoral, de la probabilidad de que el juego siga y de lo que pese la reputación. Los gobiernos que se ven en todas las cumbres tienen motivos para cooperar sin contrato; los que esperan perder las elecciones, menos.
\end{itemize}

\paragraph{c) Valoración}
\begin{itemize}
\item \emph{Por qué no basta}:
  \begin{itemize}
  \item las cumbres son pocas y los gobiernos cambian, de modo que $\beta$ es bajo;
  \item las perturbaciones y los cambios de condiciones hacen difícil saber si el otro ha cumplido o ha tenido mala suerte;
  \item el castigo, dejar de cooperar con un socio grande, es poco creíble.
  \end{itemize}
\item Por eso la cooperación tácita es frágil y pide reglas e instituciones que la vigilen (\ref{sec:formas}).
\end{itemize}

\begin{itemize}
\item Recapitulando el apartado: hay desbordamientos, aunque su signo depende del modelo (condición necesaria); el equilibrio sin coordinar es ineficiente, pero la solución cooperativa no se sostiene sola y las ganancias medidas son pequeñas (condición suficiente). La teoría explica por qué se coordina poco; la historia dirá cuándo, pese a todo, se ha coordinado.
\end{itemize}

% ==========================================================================
\section{La cooperación económica\cortetitulo en la historia macroeconómica reciente}\label{sec:historia}
\minutos{6}

\begin{itemize}
\item Una vez vista la teoría, pasamos a la historia: cuándo se ha coordinado, por qué se coordina tan poco y de qué formas.
\item \emph{Conclusión anticipada}:
  \begin{itemize}
  \item la coordinación de las políticas de demanda ha sido rara y casi siempre en crisis: Bonn en 1978, el Plaza y el Louvre en los 80, Londres en 2009;
  \item cada episodio se lee como uno de los juegos del apartado anterior, y el mismo año puede jugarse con matrices distintas: los gobiernos no discrepan tanto sobre si cooperar como sobre el signo del desbordamiento;
  \item la coordinación que dura es la de las reglas y las instituciones.
  \end{itemize}
\end{itemize}

% --------------------------------------------------------------------------
\subsection{Ejemplos de cooperación}\label{sec:ejemplos}

\begin{itemize}
\item \emph{La línea del tiempo} (tabla~\ref{tab:cronologia}): la coordinación tuvo su cima con el G-7 en los 70 y los 80, cayó en desgracia y volvió con la crisis financiera global, de la mano del G-20\nocite{Frankel2016Coordination}.
\end{itemize}

\begin{table}[H]
\centering
\caption{La coordinación internacional de las políticas económicas: hitos}\label{tab:cronologia}
\small
\begin{tabular}{>{\centering\arraybackslash}m{1.6cm}|>{\raggedright\arraybackslash}m{5.6cm}|>{\raggedright\arraybackslash}m{6.4cm}}
\rowcolor{tceeazulpalido}\textbf{Año} & \textbf{Hito} & \textbf{Lectura}\\
\hline
1944 & Bretton Woods; el plan de \autor{Keynes} (1943) no prospera & Reglas: paridades fijas pero ajustables; ajuste solo del deudor\\
\hline
1978 & Cumbre de Bonn del G-7 & \textcolor{tceeverde}{Locomotora}: estímulo de Alemania y Japón\\
\hline
1985 & Acuerdo del Plaza (G-5) & Acuerdo puntual: depreciar el dólar\\
\hline
1986-1987 & Cumbre de Tokio; acuerdo del Louvre (G-6) & Vigilancia con indicadores; estabilizar los tipos de cambio\\
\hline
1999 & G-20 de ministros y gobernadores; Foro de Estabilidad Financiera & Información: dar voz a los emergentes tras la crisis asiática\\
\hline
2008-2009 & Bajada conjunta de tipos de 6 bancos centrales; G-20 de líderes en Washington, Londres y Pittsburgh & \textcolor{tceeverde}{Locomotora} mundial; nace el Consejo de Estabilidad Financiera\\
\hline
2010 & G-20 de Toronto & \textcolor{tceedorado}{Disciplina fiscal}\\
\hline
2013 & Declaración del G-7 sobre tipos de cambio & \textcolor{tceerojo}{Guerra de divisas}: alto el fuego\\
\hline
2020-2021 & Líneas de canje de dólares; suspensión del servicio de la deuda; acuerdo de 2 pilares & Crisis de la pandemia y reglas fiscales comunes\\
\hline
2025 & Cumbre de Johannesburgo sin Estados Unidos & Foros divididos\\
\hline
\end{tabular}
\fuente{elaboración propia a partir de \textcite{Frankel2016Coordination}, \textcite{G7Bonn1978Declaration}, \textcite{G51985Plaza}, \textcite{G61987Louvre}, \textcite{G201999Berlin}, \textcite{FED2008Conjunta}, \textcite{G20Londres2009}, \textcite{G202010Toronto}, \textcite{G72013Divisas}, \textcite{Fed2020Swaps}, \textcite{G202020DSSI}, \textcite{G20Roma2021} y \textcite{G20Johannesburgo2025}.}
\end{table}

\paragraph{a) La política fiscal: de la locomotora a la disciplina}
\begin{itemize}
\item \emph{Bonn, 1978}: el primer juego de la locomotora.
  \begin{itemize}
  \item En la cumbre de julio de 1978, Alemania se comprometió a nuevas medidas de hasta el 1~\% de su PNB, Francia a elevar su déficit medio punto del PNB y Japón a crecer 1,5 puntos más; Estados Unidos, a subir el precio interior del petróleo al nivel mundial\nocite{G7Bonn1978Declaration}.
  \item Alemania se arrepintió: con la segunda crisis del petróleo, estimular resultó ser el objetivo equivocado\nocite{Frankel2016Coordination}.
  \end{itemize}
\item \emph{Londres, 2009}: la locomotora a escala mundial.
  \begin{itemize}
  \item El 2 de abril de 2009, el G-20 anunció una expansión fiscal concertada de 5 billones de dólares hasta finales de 2010 y un programa adicional de 1,1 billones, con los recursos del FMI triplicados hasta 750.000 millones\nocite{G20Londres2009} \vertema{3.A.40}.
  \item Volvamos a la matriz de la locomotora (tabla~\ref{tab:locomotora}): es el paso de (No expandir, No expandir), con 1 cada uno, a (Expandir, Expandir), con 5. Que se vieran en todas las cumbres subía el factor de descuento de \eqref{eq:gatillo}.
  \end{itemize}
\item \emph{Toronto, 2010}: el acuerdo duró poco.
  \begin{itemize}
  \item Con la crisis griega, la lectura europea fue que el déficit era el problema, y en junio de 2010 el G-20 se comprometió a reducir al menos a la mitad los déficits de las economías avanzadas antes de 2013\nocite{G202010Toronto}.
  \item En Estados Unidos, la expansión fiscal también se rechazó cuando los republicanos ganaron la Cámara de Representantes en noviembre de 2010\nocite{CamaraEEUU2026Divisiones}.
  \item Entre países, la combinación de políticas pasó de las 2 expansivas en casi todos (2009) a la consolidación fiscal en todos, con políticas monetarias que divergían: la Reserva Federal seguía comprando bonos y el BCE subió los tipos en 2011 \vertema{3.A.40}.
  \end{itemize}
\item \emph{El juego de la disciplina fiscal}\nocite{Frankel2016Coordination}:
  \begin{itemize}
  \item para Alemania y el norte de Europa, el déficit de un socio es un desbordamiento \emph{negativo}: sube el tipo de interés común o crea la expectativa de un rescate (riesgo moral);
  \item con los mismos pagos que la locomotora, la estrategia cooperativa es ahora consolidar, y la coordinación son reglas: Maastricht y el Pacto de Estabilidad y Crecimiento, reforzado en 2011\nocite{UE2011Reglamento1175} y reformado en 2024\nocite{UE2024Reglamento1263} \vertema{3.B.45};
  \item el mismo año se jugó con 2 matrices distintas: Alemania veía disciplina y Estados Unidos, locomotora. El desacuerdo no era sobre cooperar, sino sobre el signo del desbordamiento. Las tablas de este juego y del siguiente están en el anexo~\ref{anexo:juegos}.
  \end{itemize}
\end{itemize}

\paragraph{b) La política monetaria: la guerra de divisas}
\begin{itemize}
\item \emph{Los hechos}:
  \begin{itemize}
  \item la Reserva Federal hizo 3 rondas de compras de activos, en noviembre de 2008, noviembre de 2010 y septiembre de 2012\nocite{Fed2008QE,Fed2010FOMCnov,Fed2012FOMCsep}; el Banco de Japón lanzó la expansión cuantitativa y cualitativa en abril de 2013\nocite{BoJ2013QQE}, y el BCE empezó a comprar deuda pública en marzo de 2015\nocite{BCE2015APP};
  \item en cada caso la moneda se depreció y los socios se quejaron: el ministro brasileño \autor{Guido Mantega} habló de «guerra de divisas» en septiembre de 2010 y la presidenta \autor{Dilma Rousseff}, de «tsunami monetario» en 2012\nocite{Frankel2016Coordination} \vertema{3.B.15}.
  \end{itemize}
\item \emph{El juego}: otra vez un dilema del prisionero, ahora con contener como estrategia cooperativa y depreciar como dominante. Si todos deprecian, nadie gana competitividad.
\item \emph{El acuerdo}: el 12 de febrero de 2013, los ministros y gobernadores del G-7 declararon que sus políticas se orientarían a objetivos internos con instrumentos internos y que no fijarían objetivos de tipo de cambio\nocite{G72013Divisas}. Para \textcite{Frankel2016Coordination}, es el ejemplo reciente más importante de coordinación monetaria.
\item \emph{Un matiz}: con los tipos en su límite inferior, una expansión monetaria tira también de la demanda mundial. El juego se parecía más a una locomotora que a una guerra de divisas, como en los 30, cuando las devaluaciones ayudaron a salir de la depresión\nocite{EichengreenSachs1985,Frankel2016Coordination}.
\item \emph{Y en las crisis, juntos}: el 8 de octubre de 2008, 6 bancos centrales bajaron los tipos a la vez\nocite{FED2008Conjunta}, y la Reserva Federal abrió líneas de canje (\emph{swaps}) de dólares a otros bancos centrales, que amplió en marzo de 2020\nocite{Fed2020Swaps} \vertema{3.B.13}.
\end{itemize}

\paragraph{c) Los tipos de cambio: del Plaza al Louvre}
\begin{itemize}
\item En el \concepto{acuerdo del Plaza}, el 22 de septiembre de 1985, los ministros y gobernadores del G-5 declararon deseable «una nueva apreciación ordenada» de sus monedas frente al dólar y se ofrecieron a intervenir juntos\nocite{G51985Plaza}. El dólar, que ya bajaba, siguió cayendo\nocite{Frankel2015Plaza}.
\item La cumbre de Tokio, en mayo de 1986, creó la vigilancia con indicadores, y en el \concepto{acuerdo del Louvre}, el 22 de febrero de 1987, el G-6 se comprometió a «fomentar la estabilidad de los tipos de cambio en torno a los niveles actuales»\nocite{G61987Louvre}.
\item Se leen como una \concepto{batalla de los sexos} \vertema{3.A.14}: todos querían un dólar más bajo y luego estable, pero discrepaban en el nivel y en quién ajustaba su política. Hay 2 equilibrios, y coordinar es elegir uno.
\end{itemize}

\paragraph{d) ¿Quién ajusta, el deudor o el acreedor?}
\begin{itemize}
\item El país con déficit tiene que ajustar cuando se le acaba la financiación; el que tiene superávit puede no hacerlo nunca.
\item En Bretton Woods, \autor{John Maynard Keynes} propuso una Unión de Compensación que cobraba también a los acreedores\nocite{Keynes1943Proposals}; ganó el plan de \autor{Harry Dexter White} \vertema{3.B.20}.
\item Hoy lo intentan, sin sanciones, el informe anual del sector exterior del FMI, que pide a los países con superávit que impulsen su demanda interna\nocite{FMI2026ESR} \vertema{3.B.22}, la evaluación mutua del G-20 en su marco de crecimiento (\ref{sec:g20}) y, en la Unión Europea, el procedimiento de desequilibrios \vertema{3.B.45}.
\item Es la coordinación entre países con superávit y con déficit \vertema{3.B.12}; España vivió el lado deficitario de ese ajuste tras 2008 \vertema{4.A.21}. Un ejemplo, Alemania, los Países Bajos y Corea, está en el anexo~\ref{anexo:superavit}.
\end{itemize}

% --------------------------------------------------------------------------
\subsection{Obstáculos a la coordinación}\label{sec:obstaculos}

\begin{itemize}
\item Si coordinar es bueno para todos, ¿por qué se ve tan poco? Porque tiene costes, económicos y políticos.
\item \emph{Obstáculos económicos}:
  \begin{itemize}
  \item \emph{el polizón}: la solución cooperativa no es un equilibrio (\ref{sec:hamada}). Se resuelve con la repetición del juego (\ref{sec:repeticion}) o con un organismo que sancione, como el Órgano de Solución de Diferencias de la OMC \vertema{3.B.33};
  \item \emph{la coordinación puede ser contraproducente}: si los bancos centrales expanden a la vez, ninguna moneda se deprecia, y desaparece un freno a la tentación de sorprender con inflación. Los agentes lo anticipan y la inflación de equilibrio sube: coordinar solo a los bancos centrales, sin credibilidad frente al sector privado, puede empeorar a todos\nocite{Rogoff1985Counterproductive} \vertema{3.A.36};
  \item \emph{la incertidumbre}: sobre la situación de partida, sobre los objetivos y, sobre todo, sobre el modelo. Si los gobiernos no coinciden en cómo funciona la economía, coordinar con el modelo equivocado puede empeorar tanto como mejorar\nocite{FrankelRockett1988}. Es lo que pasó con la locomotora y la disciplina tras 2010 (\ref{sec:ejemplos});
  \item \emph{las preferencias distintas}: un gobierno da más peso al paro y otro a la inflación (los $\gamma_i$ de \eqref{eq:hamada}), y no hay un acuerdo que convenga igual a los 2;
  \item \emph{el reparto}: si las ganancias son desiguales, pueden hacer falta compensaciones.
  \end{itemize}
\item \emph{Obstáculos políticos}:
  \begin{itemize}
  \item la soberanía: los gobiernos no quieren ceder sus instrumentos, aunque la interdependencia ya les quita autonomía;
  \item los países no son jugadores únicos: cada gobierno negocia a la vez fuera y dentro, con sus parlamentos e intereses (exportadores frente a sectores protegidos, acreedores frente a deudores)\nocite{Putnam1988Diplomacy}, y los bancos centrales tienen mandatos nacionales\nocite{Rey2013Dilema};
  \item la economía política añade también un argumento a favor: una moneda muy depreciada por la política de un país alimenta el proteccionismo en sus socios, un coste que no recoge ningún modelo de bienestar y que justifica al menos vigilar los desalineamientos\nocite{BrozFrieden2013}.
  \end{itemize}
\end{itemize}

% --------------------------------------------------------------------------
\subsection{Formas de la coordinación}\label{sec:formas}

\begin{itemize}
\item Si, pese a los obstáculos, los países deciden coordinar, tienen que elegir cómo.
\item \emph{Reglas o acuerdos}:
  \begin{itemize}
  \item las \concepto{reglas} fijan de antemano qué hace cada uno: el patrón oro, las paridades de Bretton Woods, el Sistema Monetario Europeo \vertema{3.B.43} o las reglas fiscales de la Unión. Reducen los costes de negociar, la incertidumbre y la inconsistencia temporal, pero necesitan sanciones creíbles y pierden flexibilidad ante perturbaciones grandes, como mostró la cláusula general de escape del Pacto de Estabilidad en 2020 \vertema{3.B.45};
  \item los \concepto{acuerdos} puntuales responden a una situación: el Plaza y el Louvre, el estímulo de Londres de 2009. Son flexibles, pero duran lo que dura el interés de cada uno.
  \end{itemize}
\item \emph{Con líder o entre iguales}: un sistema con un país hegemónico, como Estados Unidos en Bretton Woods, facilita el acuerdo, porque el líder impone sus preferencias; por la misma razón, el resultado puede servir más al líder que al conjunto.
\end{itemize}

\begin{itemize}
\item Recapitulando el apartado: se ha coordinado poco y casi siempre en crisis, y los obstáculos lo explican; lo que dura son las reglas y las instituciones. Esto nos lleva a la economía política: qué hace falta para que los gobiernos cooperen y qué foros lo hacen posible.
\end{itemize}

% ==========================================================================
\section{Economía política de la coordinación:\cortetitulo mecanismos de implementación}\label{sec:economia-politica}
\minutos{8}

\begin{itemize}
\item Una vez vistos los episodios, pasamos a los mecanismos: cómo se consigue que gobiernos soberanos cooperen.
\item \emph{Conclusión anticipada}:
  \begin{itemize}
  \item un foro cambia el juego: lo repite, da información y hace visibles los incumplimientos;
  \item el G-20 es el foro de las crisis y de la coordinación de la demanda; la OCDE, el de las normas y la coordinación de la oferta; el resto (G-7, FMI, Consejo de Estabilidad Financiera, OMC, Unión Europea) sube distintos peldaños de la escalera de la cooperación.
  \end{itemize}
\end{itemize}

% --------------------------------------------------------------------------
\subsection{La problemática de teoría de juegos:\cortetitulo por qué hacen falta foros}\label{sec:problematica}

\begin{itemize}
\item \emph{El punto de partida}: la cooperación no se sostiene sola si el juego se juega una vez (\ref{sec:hamada}), y en el juego repetido solo si los gobiernos son pacientes y se pueden vigilar \eqref{eq:gatillo}.
\item \emph{Lo que hace un foro}:
  \begin{itemize}
  \item \emph{repite el juego}: las cumbres anuales y las reuniones de ministros convierten una partida suelta en una relación que no termina, y suben el factor de descuento;
  \item \emph{da información}: los informes y la evaluación por pares reducen la incertidumbre sobre la situación y sobre el modelo;
  \item \emph{hace visibles los incumplimientos}: con compromisos públicos y su seguimiento, desviarse se ve y se paga en reputación;
  \item \emph{liga asuntos}: lo que un país cede en un tema lo gana en otro.
  \end{itemize}
\item \emph{El diseño de un foro}: 3 decisiones, cada una con ventajas e inconvenientes.
  \begin{itemize}
  \item \emph{Amplitud}, cuántos asuntos entran en la negociación: más asuntos encarecen la negociación, pero abren compensaciones entre ellos. En la OMC es más fácil cerrar la agricultura si a la vez se negocian los productos industriales \vertema{3.B.33}.
  \item \emph{Profundidad}, con qué detalle se concreta cada compromiso: más detalle encarece la negociación, pero abarata la vigilancia, reduce los conflictos posteriores y da credibilidad; si el detalle es mucho, conviene prever cláusulas de flexibilidad, como la del Pacto de Estabilidad.
  \item \emph{Tamaño}: un grupo pequeño negocia más barato (el G-7); uno grande internaliza más desbordamientos (el G-20 desde 2008). Los acuerdos con muchos miembros tienden a pedir poco \vertema{3.B.31}.
  \end{itemize}
\item \emph{Los clubes}: un acuerdo se sostiene mejor si la ventaja de estar dentro es excluible. Es la propuesta de \textcite{Nordhaus2015Climate}\footnote{\autor{William Nordhaus} recibió el Premio Nobel de Economía en 2018, con \autor{Paul Romer}, «por integrar el cambio climático en el análisis macroeconómico a largo plazo»\nocite{Nobel2018Nordhaus}.} para el clima: un club que castiga con un arancel a quien no participa \vertema{3.B.31}.
\item \emph{La escalera de la cooperación} (figura~\ref{figu:escalera}) coloca cada foro en su peldaño: el intercambio de información (OCDE, FMI), la coordinación puntual (G-7, G-20) y las reglas permanentes (OMC, Unión Europea).
\end{itemize}

\begin{infografia}[H]
\centering
\caption{La escalera de la cooperación internacional y sus foros}\label{figu:escalera}
\grafico[0.75\textwidth]{escalera-cooperacion}
\fuente{elaboración propia a partir de \textcite{Bergin2008Palgrave} y \textcite{Frankel2016Coordination}.}
\end{infografia}

% --------------------------------------------------------------------------
\subsection{Los foros de coordinación: el G-20, la OCDE y otros}\label{sec:foros}

\begin{itemize}
\item Cada foro se cuenta con la misma plantilla: qué es, origen, miembros, cómo funciona, qué coordina y sus límites.
\end{itemize}

\subsubsection{El G-20}\label{sec:g20}

\begin{itemize}
\item \emph{Qué es}: el foro principal de la cooperación económica internacional, así designado por sus líderes en Pittsburgh en 2009\nocite{G20Pittsburgh2009}. Es informal: no tiene tratado, ni secretaría permanente, ni decisiones vinculantes, y decide por consenso.
\item \emph{Origen}:
  \begin{itemize}
  \item nació en 1999, tras la crisis asiática, como reunión de ministros de Finanzas y gobernadores de bancos centrales, para dar voz a las economías emergentes que el G-7 dejaba fuera; su primera reunión fue en Berlín, en diciembre de 1999\nocite{G201999Berlin};
  \item en noviembre de 2008, en Washington, celebró su primera cumbre de jefes de Estado y de Gobierno, y en Londres, en abril de 2009, sustituyó de hecho al G-7 en la gestión de la crisis\nocite{Frankel2016Coordination}.
  \end{itemize}
\item \emph{Miembros}: 19 países (Alemania, Arabia Saudí, Argentina, Australia, Brasil, Canadá, China, Corea, Estados Unidos, Francia, India, Indonesia, Italia, Japón, México, Reino Unido, Rusia, Sudáfrica y Turquía), la Unión Europea y, desde 2023, la Unión Africana\nocite{G20NuevaDelhi2023}. Reúnen en torno al 85~\% del PIB mundial, el 75~\% del comercio y 2 tercios de la población\nocite{G20Acerca}. España es invitado permanente.
\item \emph{Cómo funciona}:
  \begin{itemize}
  \item la presidencia rota cada año, apoyada por la anterior y la siguiente (la \emph{troika}): Sudáfrica en 2025; Estados Unidos desde el 1 de diciembre de 2025, con la cumbre de Miami los días 14 y 15 de diciembre de 2026\nocite{G20RG2026Miami,EEUU2026Miami}, y el Reino Unido en 2027, con la cumbre en Manchester\nocite{ReinoUnido2026Manchester};
  \item trabaja por 2 vías: la de los \emph{sherpas}, los representantes de los líderes, y la financiera, de ministros y gobernadores. El FMI, el Banco Mundial, la OCDE y el Consejo de Estabilidad Financiera le hacen el trabajo técnico.
  \end{itemize}
\item \emph{Qué ha coordinado}:
  \begin{itemize}
  \item la respuesta a la crisis: el estímulo y los recursos del FMI de Londres en 2009 (\ref{sec:ejemplos}) y la creación del Consejo de Estabilidad Financiera (\ref{sec:otros-foros});
  \item en Pittsburgh, en 2009, el Marco para un Crecimiento Fuerte, Sostenible y Equilibrado, con un proceso de evaluación mutua en el que el FMI analiza si las políticas de todos, juntas, son coherentes con los objetivos comunes\nocite{G20Pittsburgh2009};
  \item la reforma financiera: Basilea III, la resolución de los bancos sistémicos, los derivados extrabursátiles \vertema{3.B.27};
  \item la fiscalidad internacional: con la OCDE, el proyecto contra la erosión de bases y el traslado de beneficios y el acuerdo de 2 pilares (\ref{sec:ocde});
  \item la deuda de los países pobres: con la pandemia, la Iniciativa de Suspensión del Servicio de la Deuda aplazó más de 12.000 millones de dólares de pagos de casi 50 países entre mayo de 2020 y diciembre de 2021, y en noviembre de 2020 nació el Marco Común para reestructurarla, con el Club de París\nocite{G202020DSSI,G202020MarcoComun,G20Roma2021,BancoMundial2022DSSI} \vertema{3.B.29}.
  \end{itemize}
\item \emph{Límites}:
  \begin{itemize}
  \item sus compromisos no obligan y dependen de que los cumplan los gobiernos: la evaluación mutua no tiene sanción;
  \item con 20 miembros muy distintos, el consenso es difícil fuera de las crisis: la eficacia cae cuando pasa la urgencia, como de Londres (2009) a Toronto (2010);
  \item las tensiones geopolíticas lo han dividido: Estados Unidos no fue a la cumbre de Johannesburgo de 2025\nocite{G20Johannesburgo2025}.
  \end{itemize}
\end{itemize}

\subsubsection{La OCDE}\label{sec:ocde}

\begin{itemize}
\item \emph{Qué es}: la Organización para la Cooperación y el Desarrollo Económicos, un foro intergubernamental en el que sus miembros comparan sus políticas, se examinan unos a otros y fijan normas comunes. Su forma de coordinar es la \concepto{presión entre iguales} (\emph{peer pressure}): la evaluación de cada país por los demás.
\item \emph{Origen}: sucede a la Organización Europea de Cooperación Económica (OECE), creada en 1948 para repartir la ayuda del Plan Marshall. La fundó el Convenio de París, de 14 de diciembre de 1960, que firmaron 20 países, los europeos de la OECE más Estados Unidos y Canadá; España es miembro fundador\nocite{BOE1963OCDE}.
\item \emph{Miembros}: 38, con Costa Rica (2021) como la última adhesión. Tiene 8 candidatos: Argentina, Brasil, Bulgaria, Croacia, Perú y Rumanía, desde 2022, e Indonesia y Tailandia, desde 2024\nocite{OCDE2026Miembros}.
\item \emph{Órganos}:
  \begin{itemize}
  \item el \emph{Consejo}, el órgano de decisión, con un representante de cada miembro y de la Comisión Europea, presidido por el secretario general, \autor{Mathias Cormann}, en su segundo mandato desde junio de 2026\nocite{OCDE2025Cormann}; una vez al año se reúne a nivel ministerial;
  \item los \emph{comités}, por materias (política económica, comercio, fiscalidad, ayuda al desarrollo, competencia…), donde trabajan los expertos de los países;
  \item la \emph{Secretaría}, en París, que hace el análisis.
  \end{itemize}
\item \emph{Instrumentos}:
  \begin{itemize}
  \item \emph{el análisis}: el \emph{Economic Outlook}, con previsiones 2 veces al año, los estudios económicos de cada país, examinados por los demás, y \emph{Going for Growth}, que propone a cada miembro sus reformas prioritarias: coordinación de las políticas de oferta;
  \item \emph{las normas}: decisiones (vinculantes, como el Código de Liberalización de los Movimientos de Capital), recomendaciones (no vinculantes) y acuerdos, como el Convenio contra el soborno de funcionarios extranjeros o las Líneas Directrices para Empresas Multinacionales;
  \item \emph{el Consenso sobre créditos a la exportación con apoyo oficial}, un pacto de caballeros que limita la competencia en las condiciones de financiación y que la Unión Europea ha incorporado a su derecho \vertema{4.A.24};
  \item \emph{el Comité de Ayuda al Desarrollo}, que define qué es ayuda oficial al desarrollo \vertema{3.B.29};
  \item \emph{la fiscalidad}: con el G-20, el proyecto contra la erosión de bases y el traslado de beneficios (BEPS) y el acuerdo de 2 pilares del Marco Inclusivo de octubre de 2021, respaldado por los líderes del G-20 en Roma: un reparto de los beneficios de las mayores multinacionales hacia los países donde venden (pilar 1) y un impuesto mínimo efectivo del 15~\% para los grupos con ingresos de al menos 750 millones de euros (pilar 2)\nocite{G20Roma2021} \vertema{4.B.18}.
  \end{itemize}
\item \emph{Límites}: coordina sobre todo con normas y presión entre iguales, no con acuerdos de demanda; su reto es abrirse más allá del «club de países ricos» sin perder la capacidad de decidir por consenso.
\end{itemize}

\subsubsection{Otros foros}\label{sec:otros-foros}

\begin{itemize}
\item \emph{El G-7}: Estados Unidos, Japón, Alemania, Reino Unido, Francia, Italia y Canadá, con la Unión Europea. Fue el foro de la coordinación de los 70 y los 80 y lo sigue siendo de las economías avanzadas; Rusia, que lo amplió a G-8, está suspendida desde 2014. En diciembre de 2022 creó un Club del Clima, un foro abierto de cooperación sin aranceles, distinto de los clubes de \textcite{Nordhaus2015Climate}\nocite{G72022ClimateClub} \vertema{3.B.31}.
\item \emph{El Consejo de Estabilidad Financiera} (\emph{Financial Stability Board}, FSB): lo creó el G-20 en abril de 2009, como sucesor del Foro de Estabilidad Financiera de 1999, para coordinar a los reguladores y supervisores nacionales y a los organismos que fijan las normas. Reúne a los países del G-20 y a algunos más, entre ellos España; desde enero de 2013 es una asociación de derecho suizo con sede en Basilea, y lo preside \autor{Andrew Bailey} desde julio de 2025\nocite{G20Londres2009,FSB2026Historia,FSB2026Miembros,FSB2025Bailey} \vertema{3.B.27}.
\item \emph{El FMI}: el foro con tratado y con recursos.
  \begin{itemize}
  \item Vigila a cada país (las consultas del artículo IV) y al sistema en su conjunto: las Perspectivas de la Economía Mundial y el Informe del Sector Exterior, que compara la balanza de cada país con su norma. El último, de julio de 2026, concluye que los saldos excesivos crecieron en 2025, con China y Estados Unidos como principales causas, y que corregirlos a la vez da mejor resultado que hacerlo por separado\nocite{FMI2026ESR} \vertema{3.B.22}.
  \item Hace el trabajo técnico de la evaluación mutua del G-20.
  \item Desde 2009 hace con el FSB un ejercicio de alerta temprana: no intenta predecir cuándo llegará una crisis, sino identificar los riesgos extremos del sistema y lo que costarían\nocite{FMI2010EWE} \vertema{3.B.17}.
  \item Fija doctrina común: desde 2012 acepta las medidas sobre los flujos de capital en ciertas condiciones\nocite{FMI2012VisionInstitucional}, y en 2022 las admitió también como prevención cuando hay riesgo sistémico por deuda en divisas\nocite{FMI2022RevisionVI}.
  \end{itemize}
\item \emph{El Banco de Pagos Internacionales y el Comité de Basilea}: el foro de los bancos centrales y de las normas bancarias \vertema{3.B.27}.
\item \emph{La OMC}: la coordinación de la política comercial, con reglas y un órgano que resuelve las diferencias \vertema{3.B.33}.
\item \emph{El Club de París}: los acreedores oficiales que negocian juntos la reestructuración de la deuda bilateral, ahora con el Marco Común del G-20 \vertema{3.B.29}.
\item \emph{La Unión Europea}: la forma más intensa de coordinación, con reglas fiscales, el Semestre Europeo, el procedimiento de desequilibrios y el Eurogrupo; y, en el euro, una política monetaria única \vertema{3.B.16} \vertema{3.B.45}.
\item \emph{Otros}: la Red Internacional de Competencia, que coordina a las autoridades de competencia; el Foro Económico Mundial de Davos, que reúne a gobiernos y empresas sin decidir nada, y la Reunión Asia-Europa (ASEM), un diálogo entre las 2 regiones.
\end{itemize}

\begin{itemize}
\item \emph{Valoración de los foros}:
  \begin{itemize}
  \item han funcionado en las crisis (2008-2009, 2020) y en las normas (Basilea, el acuerdo fiscal de la OCDE);
  \item fuera de las crisis, la coordinación de la demanda se queda en compartir información y vigilarse: los peldaños bajos de la escalera;
  \item y la fragmentación geopolítica, con un G-20 que no consigue reunir a todos sus miembros, amenaza también esos peldaños.
  \end{itemize}
\item Recapitulando el apartado: los foros cambian el juego, lo repiten, dan información y vigilan; por eso la coordinación que existe vive en ellos, aunque suba pocos peldaños fuera de las crisis.
\end{itemize}

% ==========================================================================
\section*{Conclusión}
\minutos{1}

\begin{itemize}
\item \textbf{\emph{Recapitulación}}:
  \begin{itemize}
  \item A modo de conclusión, la coordinación es \emph{necesaria} cuando hay efectos desbordamiento, cuyo signo depende del modelo y de cómo se fijan los precios.
  \item Es \emph{deseable} porque el equilibrio sin coordinar es ineficiente, como muestra el diagrama de \autor{Hamada}.
  \item Es \emph{difícil}: la solución cooperativa no es un equilibrio, la repetición la sostiene solo con gobiernos pacientes y los obstáculos económicos y políticos la frenan.
  \item Por eso se ha coordinado sobre todo en las crisis y con reglas, y los foros suben la escalera de la cooperación peldaño a peldaño.
  \end{itemize}
\item \textbf{\emph{Relevancia}}:
  \begin{itemize}
  \item Los desequilibrios exteriores crecen, las tensiones comerciales y geopolíticas dividen los foros y la Reserva Federal sigue moviendo el ciclo financiero global: la pregunta de la coordinación es más actual que nunca.
  \end{itemize}
\item \textbf{\emph{Extensiones}}:
  \begin{itemize}
  \item el ciclo financiero global y el dilema de la política monetaria \vertema{3.B.15};
  \item la coordinación dentro de una unión monetaria y la gobernanza económica del euro \vertema{3.B.16} \vertema{3.B.45};
  \item los clubes climáticos \vertema{3.B.31} y la fiscalidad internacional con los 2 pilares \vertema{4.B.18};
  \item la deuda de los países de renta baja y el Marco Común \vertema{3.B.29}, y los derechos especiales de giro del FMI \vertema{3.B.22}.
  \end{itemize}
\item \textbf{\emph{Opinión}}:
  \begin{itemize}
  \item La coordinación es el «monstruo del lago Ness» del que hablaron \textcite{BlanchardOstryGhosh2013}: se habla de ella más de lo que se la ve. Pero aparece cuando más importa, en las crisis (2009, las líneas de canje, 2020), y las reglas que deja duran.
  \end{itemize}
\item \textbf{\emph{Idea final}}:
  \begin{itemize}
  \item «Cada uno en su casa y Dios en la de todos», decía el refrán. Poner orden en la propia casa es la primera regla, pero cuando la casa de uno ocupa media calle, la de todos también es suya.
  \end{itemize}
\end{itemize}

\clearpage
\appendix

% ==========================================================================
\section{Anexo: reforma y dólar en \autor{Corsetti} y \autor{Pesenti}}\label{anexo:corsetti-pesenti}

\begin{itemize}
\item \emph{La relación real de intercambio de los 2 países}. Para A, $\tau\equiv\big[\tfrac{P}{2}\big(\tfrac{1}{P_H}+\tfrac{1}{E\,P^*_H}\big)\big]^{-1}$; para B, con la misma lógica,
\begin{equation}
\tau_B\equiv\Big[\frac{P_B}{2}\Big(\frac{1}{P^*_F}+\frac{E}{P_F}\Big)\Big]^{-1},\label{eq:tau-b}
\end{equation}
con $P_B$ el índice de precios de consumo de B en su moneda, $P^*_F$ el precio en B de lo producido en B y $P_F$ el precio en A de lo producido en B, en moneda de A\nocite{CorsettiPesenti2009}.
  \begin{itemize}
  \item $1/\tau_B$ son las unidades de producción de B necesarias para comprar una cesta de consumo de B.
  \item Una depreciación de la moneda de A ($E\uparrow$) con PCP sube $P_F$ en la misma proporción: $E/P_F$ no cambia, pero los bienes de A se abaratan en B y $\tau_B$ sube. Con LCP, $P_F$ y $P^*_H$ están fijos en la moneda del comprador: $E/P_F$ sube y $\tau_B$ baja.
  \end{itemize}
\end{itemize}

\subsection*{Una reforma estructural en A}

\begin{itemize}
\item La reforma reduce el poder de mercado de las empresas de A: sube $\sigma$ y, con ella, el empleo natural, $\overline{\ell}=(\sigma-1)/(\sigma\chi)$, que crece con $\sigma$ (gráfico~\ref{fig:corsetti-pesenti-reforma}).
\item \emph{En A} (panel a): la NR se desplaza a la derecha, de $\overline{\ell}=3$ a $\overline{\ell}'=3{,}6$; A produce más, sus bienes se abaratan y su relación de intercambio empeora: la AS gira hacia abajo, de $C=\tfrac{2}{3}\ell$ a $C=0{,}6\,\ell$. La AD sube y A pasa de $O$ a $O'=(3{,}6;\,2{,}16)$: trabaja y consume más, aunque el empeoramiento de la relación de intercambio le quita parte de la ganancia.
\item \emph{En B} (panel b): su NR no cambia, su relación de intercambio mejora y la AS$_B$ gira hacia arriba, hasta $C=0{,}75\,\ell$; la AD$_B$ sube a 2,25 y B pasa de $O$ a $O_B'=(3;\,2{,}25)$: consume más con el mismo trabajo.
\item Las reformas de oferta de un país benefician también a sus socios: es un desbordamiento positivo\nocite{CorsettiPesenti2009}.
\end{itemize}

\begin{figure}[H]
\centering
\caption{Una reforma estructural en A ($\sigma\uparrow$) en el modelo de \autor{Corsetti} y \autor{Pesenti}}\label{fig:corsetti-pesenti-reforma}
\grafico[0.9\textwidth]{corsetti-pesenti-reforma}
\fuente{elaboración propia a partir de \textcite{CorsettiPesenti2009}.}
\end{figure}

\subsection*{La facturación en dólares}

\begin{itemize}
\item Todo el comercio se factura en la moneda de uno de los 2 países, el «dólar»: en el ejemplo, la de Estados Unidos (B). A, la zona del euro, no lo emite.
\item \emph{Si B, el emisor, expande su política monetaria}: sus precios de importación están fijos en su propia moneda, de modo que B no nota el movimiento del tipo de cambio y se comporta como una economía cerrada. Queda aislado de lo que pasa fuera.
\item \emph{En A, el no emisor} (gráfico~\ref{fig:corsetti-pesenti-dolar}):
  \begin{itemize}
  \item el dólar se deprecia y lo que A importa, facturado en dólares, se abarata en euros: baja su índice de precios y la AD sube, hasta 2,55;
  \item la AS no gira, porque las exportaciones de A también están fijas en dólares; A pasa de $O$ a $D=(3{,}825;\,2{,}55)$, sobre su AS y a la derecha de su NR, con un empleo $\ell_D$;
  \item consume como con PCP (punto $H$, en gris), pero trabaja más.
  \end{itemize}
\item \emph{El efecto neto sobre A es positivo}: por el poder de mercado, en el equilibrio inicial se trabaja menos de lo eficiente y la curva de indiferencia que pasa por $O$ corta la AS desde arriba. Moverse a lo largo de la AS hacia más empleo y más consumo sube el bienestar, siempre que el movimiento no sea demasiado grande.
\item \emph{Para la coordinación}: la política del emisor mueve el bienestar del otro país sin que el emisor lo note. Cooperar sube el bienestar mundial, pero solo tiene que cambiar su política el país que emite el dólar: la contribución es unilateral\nocite{CorsettiPesenti2009}. Es la versión de 2 países del paradigma de la moneda dominante\nocite{Gopinath2020DCP}.
\end{itemize}

\begin{figure}[H]
\centering
\caption{Una expansión monetaria del emisor del dólar vista desde el país que no lo emite}\label{fig:corsetti-pesenti-dolar}
\grafico[0.6\textwidth]{corsetti-pesenti-dolar}
\fuente{elaboración propia a partir de \textcite{CorsettiPesenti2009}.}
\end{figure}

% ==========================================================================
\section{Anexo: los juegos de la coordinación fiscal y monetaria}\label{anexo:juegos}

\begin{itemize}
\item Los 3 juegos de \textcite{Frankel2016Coordination} tienen los mismos pagos (5, 6, 0 y 1, los del dilema del prisionero de la guerra de precios \vertema{3.A.14}); solo cambian los rótulos y cuál es la estrategia cooperativa. La locomotora está en la tabla~\ref{tab:locomotora}.
\end{itemize}

\subsection*{La disciplina fiscal}

\begin{itemize}
\item Alemania (filas) y el resto de la zona del euro (columnas) eligen entre consolidar y gastar. El déficit de cada uno perjudica al otro, porque sube el tipo de interés común o porque crea la expectativa de un rescate (tabla~\ref{tab:disciplina}).
\item \emph{Sin coordinar}: gastar es una estrategia dominante, y todos acaban con déficits excesivos: la posibilidad del rescate elimina el freno.
\item \emph{Coordinando}: reglas que limiten el déficit y la deuda, como las de Maastricht, el Pacto de Estabilidad y Crecimiento de 2011 y su reforma de 2024\nocite{UE2011Reglamento1175,UE2024Reglamento1263} \vertema{3.B.45}.
\item \emph{La experiencia}: la posición alemana tras 2010 tiene 2 lecturas.
  \begin{itemize}
  \item Para unos, Alemania esperaba que gastaran los demás para mantener su superávit comercial: jugaba a la locomotora y se ponía de polizón.
  \item Para otros, jugaba de verdad al juego de la disciplina, que es el equivocado con los tipos en su límite inferior y con la demanda mundial hundida.
  \item En los 2 casos, la diferencia está en cómo se percibe el desbordamiento, no en la voluntad de cooperar\nocite{Frankel2016Coordination}.
  \end{itemize}
\end{itemize}

\begin{table}[H]
\centering
\caption{El juego de la disciplina fiscal: Alemania (filas) y el resto de la zona del euro (columnas)}\label{tab:disciplina}
\begin{tabular}{>{\centering\arraybackslash}m{2.6cm}|>{\centering\arraybackslash}m{2.6cm}|>{\centering\arraybackslash}m{2.6cm}|}
\multicolumn{1}{c}{} & \multicolumn{1}{c}{\textbf{Consolidar}} & \multicolumn{1}{c}{\textbf{Gastar}}\\
\hhline{~|-|-|}
\textbf{Consolidar} & $(5,5)$ & $(0,\underline{6})$\\
\hhline{~|-|-|}
\textbf{Gastar} & $(\underline{6},0)$ & \cellcolor{tceesalvia}$(\underline{1},\underline{1})$\\
\hhline{~|-|-|}
\end{tabular}
\fuente{elaboración propia a partir de \textcite{Frankel2016Coordination}. Subrayadas, las mejores respuestas; sombreado, el equilibrio de \autor{Nash}.}
\end{table}

\subsection*{La guerra de divisas}

\begin{itemize}
\item La zona del euro (filas) y Estados Unidos (columnas) eligen entre contener su política monetaria y expandirla para depreciar (tabla~\ref{tab:guerra-divisas}).
\item \emph{Sin coordinar}: depreciar es una estrategia dominante; los 2 deprecian, ninguno consigue la depreciación ni el estímulo exterior, y los 2 acaban con más inflación.
\item \emph{Coordinando}: abstenerse de las depreciaciones competitivas, como hizo el G-7 en febrero de 2013\nocite{G72013Divisas}.
\item \emph{El matiz}: con los tipos en el límite inferior, la expansión de cada uno tira de la demanda del otro, y el juego se vuelve una locomotora en la que coordinar es expandir juntos (\ref{sec:ejemplos}).
\end{itemize}

\begin{table}[H]
\centering
\caption{La guerra de divisas: la zona del euro (filas) y Estados Unidos (columnas)}\label{tab:guerra-divisas}
\begin{tabular}{>{\centering\arraybackslash}m{2.6cm}|>{\centering\arraybackslash}m{2.6cm}|>{\centering\arraybackslash}m{2.6cm}|}
\multicolumn{1}{c}{} & \multicolumn{1}{c}{\textbf{Contener}} & \multicolumn{1}{c}{\textbf{Depreciar}}\\
\hhline{~|-|-|}
\textbf{Contener} & $(5,5)$ & $(0,\underline{6})$\\
\hhline{~|-|-|}
\textbf{Depreciar} & $(\underline{6},0)$ & \cellcolor{tceesalvia}$(\underline{1},\underline{1})$\\
\hhline{~|-|-|}
\end{tabular}
\fuente{elaboración propia a partir de \textcite{Frankel2016Coordination}. Subrayadas, las mejores respuestas; sombreado, el equilibrio de \autor{Nash}.}
\end{table}

\subsection*{El Plaza y el Louvre: una batalla de los sexos}

\begin{itemize}
\item En una batalla de los sexos los 2 jugadores prefieren ponerse de acuerdo a no hacerlo, pero cada uno prefiere un acuerdo distinto: hay 2 equilibrios de \autor{Nash} y el problema es elegir uno \vertema{3.A.14}.
\item En 1985 y 1987, los grandes países querían frenar el dólar y luego estabilizarlo, pero discrepaban en el nivel y en quién tenía que ajustar su política. Los acuerdos del Plaza y del Louvre sirvieron para elegir un equilibrio\nocite{G51985Plaza,G61987Louvre}.
\end{itemize}

% ==========================================================================
\section{Anexo: superávit y expansión coordinada}\label{anexo:superavit}

\begin{itemize}
\item En 2019, con el comercio mundial frenándose, \textcite{GarciaAndres2019Oberturas} se preguntaba si Alemania, los Países Bajos y Corea iniciaban una expansión fiscal coordinada como la de 2009:
  \begin{itemize}
  \item los 3 tenían superávit exterior estructural: llevaban años gastando menos de lo que producían y viviendo de la demanda de otros;
  \item tenían también superávit público y pleno empleo: margen fiscal sin usar, pese a que el FMI les pedía que impulsaran la inversión;
  \item con la demanda exterior en retroceso, ese margen podía servir para sostener el crecimiento propio y el de todos.
  \end{itemize}
\item \emph{Qué ha cambiado} (gráfico~\ref{fig:saldos-superavit}):
  \begin{itemize}
  \item en 2025, los 3 siguen con superávit exterior, pero ya tienen déficit público: Alemania, del 2,9~\% del PIB; los Países Bajos, del 1,6~\%, y Corea, del 2,3~\%\nocite{FMI2026WEOdatos};
  \item Alemania ha empezado a usar su margen: en marzo de 2025 reformó su freno constitucional a la deuda para excluir parte del gasto en defensa y crear un fondo de inversión de hasta 500.000 millones de euros en 12 años\nocite{Bundestag2025Grundgesetz};
  \item el superávit exterior no ha desaparecido: el ajuste simétrico que pedía \autor{Keynes} sigue pendiente (\ref{sec:ejemplos}).
  \end{itemize}
\end{itemize}

\begin{figure}[H]
\centering
\caption{Superávit exterior y saldo público: Alemania, Países Bajos y Corea, 2019 y 2025 (\% del PIB)}\label{fig:saldos-superavit}
\grafico[0.8\textwidth]{saldos-corrientes-fiscales-superavit}
\fuente{elaboración propia a partir de \textcite{FMI2026WEOdatos} (WEO de octubre de 2026); segunda comprobación con \textcite{Eurostat2026gov10dd,Eurostat2026bopgdp6q} y \textcite{OCDE2026EO119}. En Corea, el saldo público de 2025 es estimación del FMI.}
\end{figure}

% ==========================================================================
\section{Anexo: preguntas de test}\label{anexo:test}

\begin{anotaciones}
Preguntas tipo test del primer ejercicio: \href{https://www.victorgutierrezmarcos.es/oposicion/temario/primer-ejercicio/test/simulador.html?tema=3.B.19}{simulador de test}.
\end{anotaciones}

\bibliografia

\end{document}
